% ==============================================================
% config.tex — RBRH LaTeX template
% ==============================================================
% EDITORS ONLY. Authors should never need to open this file.
%
% Owns: packages, language switch, fonts, bibliography engine,
% colors, header/footer, journal-facing translations, page
% layout, the two-column "camera-ready" paracol skeleton, title
% and author-block formatting, the author-facing figure and table
% macros, and all typographic styling (sections, captions,
% theorems, listings, admonitions, boxes, watermark, appendix
% renumbering).
%
% Ported from a real accepted RBRH submission (legacy source in
% prompts/inputs/ABRH, Collischonn et al. 2025, DOI
% 10.1590/2318-0331.302520240119), stripped of the unrelated
% business-proposal apparatus (parties, payment, signatures,
% protocols, authorization, services, team) that lived alongside
% it in the original "DocumentTeX" framework.
% ==============================================================

% \DocLanguage controls journal-facing labels (see Translations,
% below) and hyphenation. main.tex sets it before % ==============================================================
% config.tex — RBRH LaTeX template
% ==============================================================
% EDITORS ONLY. Authors should never need to open this file.
%
% Owns: packages, language switch, fonts, bibliography engine,
% colors, header/footer, journal-facing translations, page
% layout, the two-column "camera-ready" paracol skeleton, title
% and author-block formatting, the author-facing figure and table
% macros, and all typographic styling (sections, captions,
% theorems, listings, admonitions, boxes, watermark, appendix
% renumbering).
%
% Ported from a real accepted RBRH submission (legacy source in
% prompts/inputs/ABRH, Collischonn et al. 2025, DOI
% 10.1590/2318-0331.302520240119), stripped of the unrelated
% business-proposal apparatus (parties, payment, signatures,
% protocols, authorization, services, team) that lived alongside
% it in the original "DocumentTeX" framework.
% ==============================================================

% \DocLanguage controls journal-facing labels (see Translations,
% below) and hyphenation. main.tex sets it before % ==============================================================
% config.tex — RBRH LaTeX template
% ==============================================================
% EDITORS ONLY. Authors should never need to open this file.
%
% Owns: packages, language switch, fonts, bibliography engine,
% colors, header/footer, journal-facing translations, page
% layout, the two-column "camera-ready" paracol skeleton, title
% and author-block formatting, the author-facing figure and table
% macros, and all typographic styling (sections, captions,
% theorems, listings, admonitions, boxes, watermark, appendix
% renumbering).
%
% Ported from a real accepted RBRH submission (legacy source in
% prompts/inputs/ABRH, Collischonn et al. 2025, DOI
% 10.1590/2318-0331.302520240119), stripped of the unrelated
% business-proposal apparatus (parties, payment, signatures,
% protocols, authorization, services, team) that lived alongside
% it in the original "DocumentTeX" framework.
% ==============================================================

% \DocLanguage controls journal-facing labels (see Translations,
% below) and hyphenation. main.tex sets it before % ==============================================================
% config.tex — RBRH LaTeX template
% ==============================================================
% EDITORS ONLY. Authors should never need to open this file.
%
% Owns: packages, language switch, fonts, bibliography engine,
% colors, header/footer, journal-facing translations, page
% layout, the two-column "camera-ready" paracol skeleton, title
% and author-block formatting, the author-facing figure and table
% macros, and all typographic styling (sections, captions,
% theorems, listings, admonitions, boxes, watermark, appendix
% renumbering).
%
% Ported from a real accepted RBRH submission (legacy source in
% prompts/inputs/ABRH, Collischonn et al. 2025, DOI
% 10.1590/2318-0331.302520240119), stripped of the unrelated
% business-proposal apparatus (parties, payment, signatures,
% protocols, authorization, services, team) that lived alongside
% it in the original "DocumentTeX" framework.
% ==============================================================

% \DocLanguage controls journal-facing labels (see Translations,
% below) and hyphenation. main.tex sets it before \input{config};
% this \providecommand only guards against main.tex forgetting to,
% and matches the template's own default: the document is English.
\providecommand{\DocLanguage}{english}

% ==============================================================
% 1. PACKAGES
% ==============================================================

% --------------------------------------------------------------
% Fonts
% Times (serif, body text + math) + Helvetica (sans, headings/
% captions) + Source Code Pro (mono, verbatim/lstlisting).
% All bundled with standard TeX Live; compatible with pdfLaTeX.
\usepackage[T1]{fontenc}
\usepackage[utf8]{inputenc}
\usepackage{mathptmx}                    % Times, incl. matching math glyphs
\usepackage[scaled=0.92]{helvet}         % Helvetica
% Source Code Pro has a tall x-height and reads a size larger than
% Times beside it, so it is scaled down to match rather than left
% to dominate every \texttt and listing.
\usepackage[scaled=0.85]{sourcecodepro}

% --------------------------------------------------------------
% Math and symbols
\usepackage{amsmath}
\usepackage{amssymb}
\usepackage{amsthm}
\usepackage{textcomp}

% --------------------------------------------------------------
% Language, references and datetime
% cleveref and hyperref load last, after babel, per their
% required order: babel -> cleveref -> hyperref.
% Both languages are loaded, and \DocLanguage picks the MAIN one.
% "main=" is what actually decides the document language: babel
% otherwise takes the last language in the option list as main and
% re-asserts it at \begin{document}, overriding any preamble
% \selectlanguage. That is what kept captions reading "Figura" and
% the footer date in Portuguese under \DocLanguage=english.
% The option list has to be built by \edef first, because package
% options are not expanded when \usepackage reads them.
\edef\RBRHbabelopts{main=\DocLanguage,english,portuguese}
\expandafter\PassOptionsToPackage\expandafter{\RBRHbabelopts}{babel}
\usepackage{babel}
\usepackage[useregional]{datetime2}
\usepackage{csquotes}
\usepackage{hyperref}
\usepackage{cleveref}

% --------------------------------------------------------------
% Colors
\usepackage[dvipsnames, table]{xcolor}

% --------------------------------------------------------------
% Layout
\usepackage{geometry}
\usepackage{setspace}
\usepackage{titlesec}
\usepackage{fancyhdr}
\usepackage{lastpage}
\usepackage{paracol}                     % the two-column "sidebar + content" skeleton

% --------------------------------------------------------------
% Bibliography
% APA style with author-year citations, biber backend. Authors
% cite with \textcite{} in prose and \parencite{} in parentheses.
\usepackage[backend=biber, style=apa, citestyle=authoryear]{biblatex}

% --------------------------------------------------------------
% Sections
\usepackage{etoolbox}

% --------------------------------------------------------------
% Text
\usepackage{lipsum}                      % placeholder text only; not needed once authors write

% --------------------------------------------------------------
% Captions
\usepackage{caption}

% --------------------------------------------------------------
% Figures
\usepackage{graphicx}

% --------------------------------------------------------------
% Tables
\usepackage{array}
\usepackage{booktabs}
\usepackage{threeparttable}
% tabularx so a table can be told to occupy exactly the width it
% is given. threeparttable sizes the caption and the notes to the
% width of the tabular, so a table that fills \linewidth is what
% keeps both flush with the text column instead of ending short.
\usepackage{tabularx}
% Left-aligned flexible column: the X column, minus the centring.
\newcolumntype{L}{>{\raggedright\arraybackslash}X}

% --------------------------------------------------------------
% Boxes
\usepackage{tcolorbox}
\tcbuselibrary{skins, breakable}

% --------------------------------------------------------------
% Listings
\usepackage{listings}

% --------------------------------------------------------------
% Misc
\usepackage{fontawesome5}                % admonition icons
% ORCID mark for the author list. Draws the official logo as
% vector paths rather than embedding an SVG, which pdfLaTeX
% cannot do without --shell-escape. It scales the mark to the
% current font's X-height, so it sits inside the line and cannot
% change the leading, and it makes the icon itself the hyperlink.
\usepackage{orcidlink}

% ==============================================================
% 2. LANGUAGE SWITCH
% ==============================================================
% Nothing to do here: the document language is set by babel's
% "main=" option in section 1, which is the only thing that
% survives \begin{document}. It drives hyphenation, babel-aware
% dates (datetime2 [useregional]), quoting (csquotes) and the
% caption words ("Figure"/"Figura", "Table"/"Tabela").
%
% A preamble \selectlanguage used to sit here and did nothing:
% babel overrode it on entering the document body.

% ==============================================================
% 3. TRANSLATIONS
% ==============================================================
% Journal-facing labels, in Portuguese or English depending on
% \DocLanguage. Authors reference these by name (e.g. \trAbstract)
% rather than hardcoding either language, so the whole document
% switches with one macro.
% \ifx rather than etoolbox's \ifdefstring: \ifx is expandable, so
% \trlang survives \MakeUppercase. With \ifdefstring the conditional
% gets mangled during uppercasing and the wrong branch wins -- that
% is what made the sidebar print "ESTUDO DE CASO" under english.
%
% Both sides are normalised into plain macros first. \ifx compares
% meanings, and \newcommand makes \DocLanguage \long while a bare
% \def would not -- comparing them directly fails on that prefix
% alone and silently selects Portuguese. \edef flattens whatever
% the author wrote into the same shape as \RBRHenglish.
%
% No "@" in these names: this is a .tex file, not a .sty, so @ is
% not a letter here and \RBRHxx@yy would parse as \RBRHxx followed
% by stray characters.
\edef\RBRHlang{\DocLanguage}
\def\RBRHenglish{english}
\newcommand{\trlang}[2]{\ifx\RBRHlang\RBRHenglish #1\else #2\fi}

% --- Front matter ---
\newcommand{\trAbstract}{\trlang{Abstract}{Resumo}}
\newcommand{\trKeywords}{\trlang{Keywords}{Palavras-chave}}
\newcommand{\trListOfBoxes}{\trlang{Boxes}{Destaques}}
\newcommand{\trHighlights}{\trlang{Highlights}{Destaques do estudo}}
\newcommand{\trCorrespondence}{\trlang{Correspondence to}{Correspond\^encia para}}
\newcommand{\trCitation}{\trlang{Citation}{Como citar}}

% --- Cover sidebar ---
\newcommand{\trEditors}{\trlang{Editors}{Editores}}
\newcommand{\trEditorInChief}{\trlang{Editor-in-Chief}{Editor-chefe}}
\newcommand{\trAssociatedEditor}{\trlang{Associated Editor}{Editor associado}}
\newcommand{\trCompetingInterests}{\trlang{Competing interests}{Conflito de interesses}}
\newcommand{\trDataAvailability}{\trlang{Data availability}{Disponibilidade dos dados}}
\newcommand{\trPublicationDates}{\trlang{Publication dates}{Datas de publica\c{c}\~ao}}
\newcommand{\trPublished}{\trlang{Published}{Publicado em}}
\newcommand{\trIssue}{\trlang{Issue}{Edi\c{c}\~ao}}
\newcommand{\trHistory}{\trlang{History}{Hist\'orico}}
\newcommand{\trReceived}{\trlang{Received}{Recebido em}}
\newcommand{\trReviewed}{\trlang{Reviewed}{Revisado em}}
\newcommand{\trAccepted}{\trlang{Accepted}{Aceito em}}

% --- Back matter ---
\newcommand{\trReferences}{\trlang{References}{Refer\^encias}}
\newcommand{\trAppendix}{\trlang{Appendix}{Ap\^endice}}
\newcommand{\trAppendices}{\trlang{Appendices}{Ap\^endices}}
\newcommand{\trAcknowledgements}{\trlang{Acknowledgements}{Agradecimentos}}

% --- Items ---
\newcommand{\trListing}{\trlang{Code}{C\'odigo}}
\newcommand{\trListings}{\trlang{Code}{C\'odigos}}
\newcommand{\trDraft}{\trlang{DRAFT}{MINUTA}}

% --- Admonitions ---
\newcommand{\trAdmNote}{\trlang{Note}{Nota}}
\newcommand{\trAdmTip}{\trlang{Tip}{Dica}}
\newcommand{\trAdmImportant}{\trlang{Important}{Importante}}
\newcommand{\trAdmWarning}{\trlang{Warning}{Alerta}}
\newcommand{\trAdmCaution}{\trlang{Caution}{Cuidado}}

% ==============================================================
% 4. COLORS
% ==============================================================
% Override this section to apply a different brand/theme.

% --- Grayscale utilities ---
\definecolor{verylightgray}{gray}{0.97}
\definecolor{lightgray}{gray}{0.85}
\definecolor{midgray}{gray}{0.50}
\definecolor{heavygray}{gray}{0.3}

% --- Brand / accent ---
% The journal's blue is the accent colour: it carries the article
% title, the section headings, the abstract headings and the
% metadata column. The masthead itself is set in black, so the
% blue marks structure rather than branding.
%
% Everything already keyed to AccentColor follows from here --
% caption labels, theorem heads and the tcolorbox palette, whose
% tints derive from it.
\definecolor{AccentColor}{RGB}{38,52,139}
\colorlet{AccentColorVeryLight}{AccentColor!5}
\colorlet{AccentColorLight}{AccentColor!20}

% --- Semantic aliases ---
\definecolor{HeadingColor}{named}{AccentColor}
\definecolor{LinkColor}{RGB}{38,52,139}
\definecolor{RuleColor}{named}{midgray}
\definecolor{TextColor}{named}{black}

% --- Page and text ---
\pagecolor{white}
\color{TextColor}

% --- Links ---
\hypersetup{
    colorlinks=true,
    linkcolor=LinkColor,
    citecolor=LinkColor,
    filecolor=LinkColor,
    urlcolor=LinkColor,
}

% ==============================================================
% 5. HEADER AND FOOTER
% ==============================================================
% Every field comes from cover.tex, which build_cover.py generates
% from the TOML files. Both journal logos sit in the running
% header at the same height, one at each margin; the footer
% carries the short citation on the left and the page count on
% the right.
\newcommand{\RBRHlogoheight}{8mm}
\newcommand{\HeaderLeft}{\includegraphics[height=\RBRHlogoheight]{\JournalLogoLeft}}
\newcommand{\HeaderRight}{\includegraphics[height=\RBRHlogoheight]{\JournalLogoRight}}
\newcommand{\FooterLeft}{\ArticleShortCitation}
\newcommand{\FooterCenter}{}
\newcommand{\FooterRight}{\thepage\;/\;\pageref{LastPage}}

\fancyhfoffset[L]{0mm}
\fancyhfoffset[R]{0mm}

\fancypagestyle{plain}{%
  \fancyhf{}
  \fancyhead[L]{\HeaderLeft}
  \fancyhead[R]{\HeaderRight}
  \renewcommand{\headrulewidth}{0.4pt}
  \fancyfoot[L]{\sffamily\scriptsize{\FooterLeft}}
  \fancyfoot[C]{\sffamily\scriptsize{\FooterCenter}}
  \fancyfoot[R]{\sffamily\scriptsize{\FooterRight}}
  \renewcommand{\footrulewidth}{0.2pt}
}

\pagestyle{fancy}
\fancyhf{}
\fancyhead[L]{\HeaderLeft}
\fancyhead[R]{\HeaderRight}
\renewcommand{\headrulewidth}{0.4pt}
\fancyfoot[L]{\sffamily\scriptsize{\FooterLeft}}
\fancyfoot[C]{\sffamily\scriptsize{\FooterCenter}}
\fancyfoot[R]{\sffamily\scriptsize{\FooterRight}}
\renewcommand{\footrulewidth}{0.2pt}

% ==============================================================
% 6. PAGE LAYOUT AND SPACING
% ==============================================================
\geometry{
 a4paper,
 left=20mm,
 right=20mm,
 top=20mm,
 includefoot,
 bottom=20mm,
 headheight=15mm,
 headsep=5mm,
 footskip=10mm,
}

\newcommand{\pagesectionvspace}{\vspace*{0mm}}

\setlength{\parindent}{0mm}
\setlength{\parskip}{2mm}

\setlength{\intextsep}{1mm}
\setlength{\textfloatsep}{2mm}
\setlength{\floatsep}{2mm}
\setlength{\abovecaptionskip}{1mm}
\setlength{\belowcaptionskip}{0mm}

\linespread{1}

\titlespacing*{\section}      {0mm}{4mm}{0mm}
\titlespacing*{\subsection}   {0mm}{2mm}{0mm}
\titlespacing*{\subsubsection}{0mm}{2mm}{0mm}
\titlespacing{\paragraph}{0mm}{2mm}{2mm}

% --------------------------------------------------------------
% Camera-ready sidebar column (left, narrow) + content column
% (right, wide). Width tuned for the DOI/highlights/correspondence
% sidebar used by main.tex; adjust \RBRHColumnRatio to change it.
\newcommand{\RBRHColumnRatio}{0.258}

% Lengths for the small-with-side-caption figure pattern (see
% contents.tex): a small, fixed-size image whose caption spans
% into the space the narrow sidebar column would otherwise
% occupy. Allocated once here; reused via \setlength (not
% \newlength) at the use site.
\newlength{\sidecapimgwidth}
\newlength{\sidecapgap}
\newlength{\sidecapwidth}
\newlength{\smallfigfixedwidth}

% ==============================================================
% 7. BIBLIOGRAPHY RESOURCES
% ==============================================================
\addbibresource{references.bib}
% Reference list one size down from the body text.
\renewcommand*{\bibfont}{\small}

% ==============================================================
% 8. CAPTIONS
% ==============================================================
\DeclareCaptionLabelSeparator{pipe}{\;|\;}

\captionsetup[figure]{
    font={footnotesize, sf, color=heavygray},
    labelfont={bf, color=AccentColor},
    labelsep=quad,
}

% No "width=" here. Inside paracol \textwidth is still the full
% page width, not the width of the column being typeset, so
% width=\textwidth made table captions overrun the text column and
% spill into the sidebar. Left unset, the caption takes the width
% of whatever contains it -- the text column for TableM, the full
% span for TableL -- which is right in both cases.
\captionsetup[table]{
    font={footnotesize, sf},
    labelfont={bf, color=AccentColor},
    labelsep=pipe,
}

\DeclareCaptionStyle{lstlisting}{
    font={footnotesize, sf},
    labelfont={bf, color=AccentColor},
    labelsep=pipe,
    justification=centering
}
\captionsetup[lstlisting]{style=lstlisting}

% All captions: justified, flush left, regardless of length.
% singlelinecheck=false disables caption's default of auto-centering
% any caption that happens to fit on a single line.
\captionsetup{justification=justified,singlelinecheck=false,labelfont=bf}

% ==============================================================
% 9. SECTION TITLES
% ==============================================================
\titleformat{\section}
  {\bfseries\sffamily\Large\color{HeadingColor}}
  {\makebox[8mm][l]{\thesection}}{2mm}{}

\titleformat{\subsection}
  {\bfseries\sffamily\large\color{HeadingColor}}
  {\makebox[8mm][l]{\thesubsection}}{2mm}{}

\titleformat{\subsubsection}
  {\sffamily\normalsize\color{HeadingColor}}
  {\makebox[10mm][l]{\thesubsubsection}}{4mm}{}

\titleformat{\paragraph}[runin]
  {\bfseries\sffamily\normalsize\color{HeadingColor}}
  {}{0mm}{}[.]

% Shared heading for auxiliary lists (References, List of Boxes,
% Glossary, ...) so they match the plain \section* styling used
% by \tableofcontents / \listoffigures / \listoftables, instead of
% falling back to each package's own internal heading.
\newcommand{\auxlisttitle}[1]{%
  \begingroup
    \titlespacing*{\section}{0mm}{6mm}{2mm}%
    \section*{#1}%
  \endgroup
  \phantomsection
  \addcontentsline{toc}{section}{#1}%
}

% ==============================================================
% 10. THEOREM-LIKE ENVIRONMENTS
% ==============================================================
\providecommand{\theoremname}{\trlang{Theorem}{Teorema}}
\providecommand{\lemmaname}{\trlang{Lemma}{Lema}}
\providecommand{\corollaryname}{\trlang{Corollary}{Corol\'ario}}
\providecommand{\definitionname}{\trlang{Definition}{Defini\c{c}\~ao}}
\providecommand{\examplename}{\trlang{Example}{Exemplo}}
\providecommand{\remarkname}{\trlang{Remark}{Observa\c{c}\~ao}}

\makeatletter
\renewcommand{\th@plain}{%
    \thm@headfont{\bfseries\sffamily\color{AccentColor}}%
    \itshape
}
\renewcommand{\th@definition}{%
    \thm@headfont{\bfseries\sffamily\color{AccentColor}}%
    \normalfont
}
\renewcommand{\th@remark}{%
    \thm@headfont{\itshape\sffamily\color{AccentColor}}%
    \normalfont
}
\makeatother

\newtheorem{theorem}{\theoremname}
\newtheorem{lemma}[theorem]{\lemmaname}
\newtheorem{corollary}[theorem]{\corollaryname}
\newtheorem{definition}{\definitionname}
\newtheorem{example}{\examplename}
\newtheorem*{remark}{\remarkname}

% ==============================================================
% 11. LISTINGS
% ==============================================================
\renewcommand{\lstlistingname}{\trListing}
\crefname{lstlisting}{\trListing}{\trListings}
\Crefname{lstlisting}{\trListing}{\trListings}

\lstdefinestyle{codeblock}{
    backgroundcolor=\color{black!5},
    commentstyle=\color{black!40}\itshape,
    keywordstyle=\color{purple}\bfseries,
    stringstyle=\color{teal},
    numberstyle=\tiny\color{black!40},
    basicstyle=\ttfamily\footnotesize,
    breaklines=true,
    breakatwhitespace=true,
    postbreak=\mbox{\textcolor{gray}{$\hookrightarrow$}\space},
    numbersep=8pt,
    tabsize=4,
    keepspaces=true,
    frame=none,
    captionpos=b,
    showspaces=false,
    showstringspaces=false,
    showtabs=false,
    numbers=left,
}
\lstset{style=codeblock, language=Python}
\lstset{basicstyle=\ttfamily\footnotesize}
\renewcommand{\ttdefault}{SourceCodePro-TLF}
\renewcommand{\familydefault}{\rmdefault}

% ==============================================================
% 12. BOXES
% ==============================================================
\tcbset{
    boxbasestyle/.style={
        enhanced,
        breakable,
        colback      = AccentColorVeryLight,
        colbacktitle = AccentColorLight,
        coltitle     = AccentColor,
        colframe     = AccentColor,
        boxrule      = 0.5pt,
        arc          = 2pt,
        fonttitle    = \bfseries\sffamily\small,
        fontupper    = \sffamily\footnotesize,
        boxsep       = 0mm,
        top          = 2mm,
        bottom       = 2mm,
        left         = 2mm,
        right        = 2mm,
        toptitle     = 1mm,
        bottomtitle  = 1mm,
    }
}

% boxstory — numbered, breakable, document-wide counter.
% Usage: \begin{boxstory}{Title} ... \end{boxstory}
\newtcolorbox[
    auto counter,
    list inside = lob,
]{boxstory}[2][]{
    boxbasestyle,
    title = {\sffamily Box~\thetcbcounter\quad #2},
    #1,
}

\newcommand{\notitleheading}[1]{}
\newcommand{\listofboxes}{%
    \auxlisttitle{\trListOfBoxes}%
    \tcblistof[\notitleheading]{lob}{\trListOfBoxes}%
}

% ==============================================================
% 13. ADMONITIONS
% ==============================================================
% Color-coded unnumbered admonition boxes. Icons require fontawesome5.
\newtcolorbox{amdnote}{
    boxbasestyle,
    colback = cyan!5, colbacktitle = cyan!50!black, colframe = cyan!50!black,
    coltitle = white, title = {\faInfoCircle\quad \trAdmNote},
}
\newtcolorbox{amdtip}{
    boxbasestyle,
    colback = green!5!white, colbacktitle = green!40!black, colframe = green!40!black,
    coltitle = white, title = {\faLightbulb\quad \trAdmTip},
}
\newtcolorbox{amdimportant}{
    boxbasestyle,
    colback = violet!5, colbacktitle = violet!60!black, colframe = violet!60!black,
    coltitle = white, title = {\faExclamationCircle\quad \trAdmImportant},
}
\newtcolorbox{amdwarning}{
    boxbasestyle,
    colback = orange!5, colbacktitle = orange!70!black, colframe = orange!70!black,
    coltitle = white, title = {\faExclamationTriangle\quad \trAdmWarning},
}
\newtcolorbox{amdcaution}{
    boxbasestyle,
    colback = red!5, colbacktitle = red!70!black, colframe = red!70!black,
    coltitle = white, title = {\faBan\quad \trAdmCaution},
}

% ==============================================================
% 14. DRAFT WATERMARK
% ==============================================================
% main.tex sets \IsDraft to "true" to stamp every page "DRAFT" /
% "MINUTA" while the manuscript is a work in progress.
\providecommand{\IsDraft}{false}
\ifdefstring{\IsDraft}{true}{
    \usepackage{draftwatermark}
    \SetWatermarkText{\trDraft}
    \SetWatermarkScale{1}
    \SetWatermarkColor[gray]{0.8}
    \SetWatermarkAngle{30}
}{}

% ==============================================================
% 15. COVER
% ==============================================================
% Everything here reads macros defined in cover.tex, which
% build_cover.py generates from journal.toml, article.toml and
% authors.toml. Nothing on the cover is written by hand, which is
% why contents.tex can begin straight at the first section.
%
% \PrintMasthead fills the wide column: journal name, title,
% authors, affiliations, then both abstracts. \PrintSidebar fills
% the narrow one, in the order RBRH prints it.

% The author and affiliation blocks are long runs of proper nouns
% and institution names, which offer TeX very few legal break
% points. On the default settings the author list overflowed the
% column by about 10pt.
%
% Justified was tried first and does fix the overflow, with
% \tolerance raised and \emergencystretch added -- but it pays for
% it by hyphenating personal names across lines ("An-derson",
% "Fer-nando") and by leaving one very loose line. Breaking a
% surname is not acceptable in a masthead, so these blocks are set
% ragged right instead, with hyphenation switched off outright.
% The right edge is uneven; nothing overflows and no name breaks.
%
% To go back to justified, swap \raggedright and the two penalties
% for: \tolerance=4000 \emergencystretch=2em
%
% These are paragraph-breaking parameters, so they only take
% effect if the paragraph ends while they are still in scope --
% hence the \par inside each group below.
\newcommand{\RBRHnameblock}{%
  \raggedright
  \hyphenpenalty=10000
  \exhyphenpenalty=10000
}

% --- Wide column ---
\newcommand{\PrintMasthead}{%
  {\noindent \sffamily \Large\textbf{\JournalName}}\\
  {\noindent \large \textit{\JournalNameEn}}\\[3mm]
  \begin{center}
    \raggedright {\Large \sffamily \textbf{\color{AccentColor}\ArticleTitle}}\\[3mm]
  \end{center}
  {\RBRHnameblock\noindent\AuthorBlock\par}
  \vspace{2mm}
  % Affiliations in the same grey as the captions, so they read as
  % secondary to the author names without fading out like the
  % licence notice does.
  {\footnotesize\color{heavygray}\RBRHnameblock\noindent\AffiliationBlock\par}
  \vspace{2mm}
  \hrule
  \vspace{2mm}
  \PrintAbstracts
}

% Both languages are printed whatever \DocLanguage says: RBRH
% carries an English abstract and a Portuguese resumo regardless
% of the language the body is written in.
% The \small applies to the heading words only -- "Abstract",
% "Keywords:" and their Portuguese counterparts, which were
% sitting at full body size. The abstract text itself stays at
% \footnotesize.
\newcommand{\PrintAbstracts}{%
  \begin{spacing}{0.8}
  \small
  {\noindent \sffamily\textbf{\color{AccentColor}Abstract}\quad{\footnotesize\color{TextColor} \AbstractEn}}\\[1mm]
  {\noindent \sffamily\textbf{\color{AccentColor}Keywords:}\quad{\footnotesize\color{TextColor} \KeywordsEn}}\\[3mm]
  {\noindent \sffamily\textbf{\color{AccentColor}Resumo}\quad{\footnotesize\color{TextColor} \AbstractPt}}\\[1mm]
  {\noindent \sffamily\textbf{\color{AccentColor}Palavras-chave:}\quad{\footnotesize\color{TextColor} \KeywordsPt}}
  \end{spacing}
}

% --- Narrow column ---
% One helper per block, so the running order is readable and can
% be reshuffled by editing \PrintSidebar alone.
% The accent marks the structure of the column -- the heading of
% each block -- while the content under it stays in body colour.
\newcommand{\RBRHcoverblock}[2]{%
  {\color{AccentColor}\textbf{#1}}\\[1mm]
  #2
  \\[4mm]
}

\newcommand{\PrintSidebar}{%
  {
  \sffamily
  \scriptsize
  \raggedright
  % licence block at the foot overrides it with grey.
  \textbf{\MakeUppercase{\ArticleType}}\\
  \ifdefempty{\ArticleDOI}{}{DOI\;\ArticleDOI\\[4mm]}
  \RBRHcoverblock{\trHighlights:}{\ArticleHighlights}
  \RBRHcoverblock{\trCitation:}{\ArticleCitation}
  \RBRHcoverblock{\trCorrespondence:}{\CorrespondingBlock}
  \RBRHcoverblock{\trCompetingInterests:}{\CompetingInterests}
  \RBRHcoverblock{\trAcknowledgements:}{\Acknowledgements}
  \RBRHcoverblock{\trEditors:}{%
    \trEditorInChief:\\ \EditorInChief\\[1mm]
    \trAssociatedEditor:\\ \AssociatedEditor}
  \RBRHcoverblock{\trDataAvailability:}{\DataAvailability}
  \RBRHcoverblock{\trPublicationDates:}{%
    \trPublished\ \DatePublished\\
    \trIssue\ \ArticleIssue}
  \RBRHcoverblock{\trHistory:}{%
    \trReceived\ \DateReceived\\
    \trReviewed\ \DateReviewed\\
    \trAccepted\ \DateAccepted}
  % Push the licence to the foot of the column: everything above
  % sets at its natural height and this absorbs the slack, so the
  % notice sits on the baseline of the page however much or little
  % metadata precedes it.
  \vfill
  \PrintLicense
  }
}

% CC-BY notice, set in grey so it reads as a footer to the
% sidebar rather than as content. The badge is a vector PDF
% converted from assets/cc-by.svg.
\newcommand{\PrintLicense}{%
  {\color{midgray}
   \includegraphics[height=3.2mm]{\LicenseBadge}\\[1mm]
   \LicenseCopyright\\
   \LicenseStatement\par}
}

% ==============================================================
% 16. APPENDIX RENUMBERING
% ==============================================================
% Call \BeginAppendix in contents.tex to switch into lettered,
% restarted numbering (A, A.1, Figure A.1, ...) for supplementary
% material.
\newcommand{\BeginAppendix}{%
  \appendix
  \renewcommand{\thesection}{\Alph{section}}
  \renewcommand{\thesubsection}{\Alph{section}.\arabic{subsection}}
  \renewcommand{\thesubsubsection}{\Alph{section}.\arabic{subsection}.\arabic{subsubsection}}
  \setcounter{section}{0}
  \setcounter{equation}{0}
  \setcounter{figure}{0}
  \setcounter{table}{0}
  \setcounter{lstlisting}{0}
  \renewcommand{\theequation}{\Alph{section}.\arabic{equation}}
  \renewcommand{\thefigure}{\Alph{section}.\arabic{figure}}
  \renewcommand{\thetable}{\Alph{section}.\arabic{table}}
  \renewcommand{\thelstlisting}{\Alph{section}.\arabic{lstlisting}}
  \pagesectionvspace
  \auxlisttitle{\trAppendices}
  \label{sec:appendix}
  \vspace{10mm}
  \small
  % Headings one level down too. The body is already \small, which
  % left \Large section titles looking oversized against it.
  \titleformat{\section}
    {\bfseries\sffamily\large\color{HeadingColor}}
    {\makebox[8mm][l]{\thesection}}{2mm}{}
  \titleformat{\subsection}
    {\bfseries\sffamily\normalsize\color{HeadingColor}}
    {\makebox[8mm][l]{\thesubsection}}{2mm}{}
  \titleformat{\subsubsection}
    {\sffamily\small\color{HeadingColor}}
    {\makebox[10mm][l]{\thesubsubsection}}{4mm}{}
}

% ==============================================================
% 17. SHARED COUNTERS ACROSS PARACOL COLUMNS
% ==============================================================
% paracol makes every counter column-LOCAL by default -- its
% internal list of global counters ships as just {page}. Each
% column therefore keeps its own `figure`, `table` and `equation`.
% A figure written inside a column and one written in a
% \switchcolumn[1]*[...] spanning block are stepped in different
% counter states, so the spanning one restarts at 1 and the
% numbering appears to repeat. `figure` is shared instead: the two
% columns here are a metadata sidebar plus one article body, not
% parallel independent texts.
%
% Name the counter explicitly -- do NOT use \globalcounter*, and
% do NOT add `table` here. Both were tried and both break this
% document with "Missing \endcsname inserted" at the first
% \section (the starred form because it sweeps in every counter,
% including `table`). Bisected against this exact package set:
%   figure   -> safe, and is what fixes the numbering bug
%   equation -> safe, but unused; add only if an equation ever
%               goes inside a \switchcolumn[1]*[...] span
%   table    -> BREAKS the build, do not enable
% `footnote` is left alone too: paracol manages its scope itself.
\globalcounter{figure}

% ==============================================================
% 18. AUTHOR-FACING FIGURE AND TABLE MACROS
% ==============================================================
% So contents.tex never carries layout plumbing:
%
%   \FigureM[pos]{name}{caption}{label}        text width
%   \FigureS[pos][frac]{name}{caption}{label}  small -- frac of the
%                                              text width, default
%                                              0.67 -- and the only
%                                              size whose caption
%                                              sits beside the
%                                              image, spanning into
%                                              the sidebar column
%   \FigureL[pos]{name}{caption}{label}        full width, margin to
%                                              margin, matching the
%                                              width of the header
%                                              rule
%
%   \begin{TableM}[pos]{caption}{label} ... \end{TableM}  text width
%   \begin{TableL}[pos]{caption}{label} ... \end{TableL}  full width
%
% [pos] is the usual LaTeX float placement specifier and defaults
% to t (top). The table environments take the tabular material as
% their body and put the caption above it, per table convention;
% figures caption below.
%
% <name> is the bare file name: no folder, no extension. The file
% is resolved inside figures/, preferring .pdf, then .png, then
% .jpeg, then .jpg. A name that matches nothing degrades to a
% placeholder box rather than failing the build.
%
% M/S/L map onto real floats -- figure/table for M, and the
% double-column figure*/table* for S and L, which is what makes
% them span the sidebar. paracol supports both (it carries its own
% \pcol@floatplacement and spanning-float machinery), and the
% figure counter is shared across columns by section 19 above, so
% numbering follows source order. \label goes straight after
% \caption inside the float, where \@currentlabel is correct.

\newif\ifRBRHfigfound
\newcommand{\RBRHfigpath}{}

% Resolve figures/<name> against the supported extensions, in
% order of preference.
\newcommand{\RBRHfindfigure}[1]{%
  \IfFileExists{figures/#1.pdf}%
    {\RBRHfigfoundtrue \def\RBRHfigpath{figures/#1.pdf}}%
    {\IfFileExists{figures/#1.png}%
      {\RBRHfigfoundtrue \def\RBRHfigpath{figures/#1.png}}%
      {\IfFileExists{figures/#1.jpeg}%
        {\RBRHfigfoundtrue \def\RBRHfigpath{figures/#1.jpeg}}%
        {\IfFileExists{figures/#1.jpg}%
          {\RBRHfigfoundtrue \def\RBRHfigpath{figures/#1.jpg}}%
          {\RBRHfigfoundfalse \def\RBRHfigpath{}}}}}%
}

% Placeholder for a missing file: gray box in golden ratio
% (height = width / 1.618) carrying a red notice. \detokenize
% keeps names containing "_" printable.
\newlength{\RBRHfigwidth}
\newlength{\RBRHfigheight}
\newcommand{\RBRHmissingfigure}[2]{%
  \setlength{\RBRHfigwidth}{#1}%
  \setlength{\RBRHfigheight}{\dimexpr\RBRHfigwidth*1000/1618\relax}%
  \begingroup
    \setlength{\fboxsep}{0pt}%
    \colorbox{lightgray}{%
      \begin{minipage}[c][\RBRHfigheight][c]{\RBRHfigwidth}%
        \centering
        \textcolor{red}{\sffamily\small
          Figure \texttt{\detokenize{#2}} not found}%
      \end{minipage}%
    }%
  \endgroup
}

% The image at #1 wide, or the placeholder if it is missing.
\newcommand{\RBRHfigbody}[2]{%
  \ifRBRHfigfound
    \includegraphics[width=#1]{\RBRHfigpath}%
  \else
    \RBRHmissingfigure{#1}{#2}%
  \fi
}

% --- M: text width ---
\NewDocumentCommand{\FigureM}{O{t} m m m}{%
  \begin{figure}[#1]
    \centering
    \RBRHfindfigure{#2}%
    \RBRHfigbody{\linewidth}{#2}%
    \caption{#3}\label{#4}%
  \end{figure}%
}

% --- L: full width, margin to margin ---
\NewDocumentCommand{\FigureL}{O{t} m m m}{%
  \begin{figure*}[#1]
    \centering
    \RBRHfindfigure{#2}%
    \RBRHfigbody{\textwidth}{#2}%
    \caption{#3}\label{#4}%
  \end{figure*}%
}

% --- S: small, caption beside the image, spanning the sidebar ---
% The image width is captured BEFORE the float opens, while
% \linewidth still means the text column: inside a double-column
% float LaTeX widens it to \textwidth. The float body is typeset
% here at the call site, so the captured value is the right one.
\NewDocumentCommand{\FigureS}{O{t} O{0.67} m m m}{%
  \setlength{\smallfigfixedwidth}{#2\linewidth}%
  \begin{figure*}[#1]
    \RBRHfindfigure{#3}%
    \setlength{\sidecapimgwidth}{\smallfigfixedwidth}%
    \setlength{\sidecapgap}{6mm}%
    \setlength{\sidecapwidth}{\textwidth}%
    \addtolength{\sidecapwidth}{-\sidecapimgwidth}%
    \addtolength{\sidecapwidth}{-\sidecapgap}%
    \noindent
    % \vspace{0pt} as the first thing in each minipage is what makes
    % [t] align the TOPS. Without it, [t] aligns first baselines, and
    % a minipage holding just a graphic has its baseline at the
    % graphic's bottom -- which dropped the caption to the foot of
    % the image instead of level with its top. \abovecaptionskip is
    % zeroed for the same reason: it is meant to separate a figure
    % from a caption printed below it, not beside it.
    \begin{minipage}[t]{\sidecapwidth}
      \vspace{0pt}%
      \setlength{\abovecaptionskip}{0pt}%
      \caption{#4}\label{#5}%
    \end{minipage}%
    \hspace{\sidecapgap}%
    \begin{minipage}[t]{\sidecapimgwidth}
      \vspace{0pt}%
      \RBRHfigbody{\linewidth}{#3}%
    \end{minipage}%
  \end{figure*}%
}

% --- Tables: caption above the body, per table convention ---
% The body is wrapped in threeparttable, so the caption and any
% notes take the width of the tabular rather than running wider,
% and authors can add a notes block without extra scaffolding:
%
%   \begin{TableM}{caption}{label}
%     \begin{tabular}{ll} ... \end{tabular}
%     \begin{tablenotes}
%       \item Note text.
%     \end{tablenotes}
%   \end{TableM}
\NewDocumentEnvironment{TableM}{O{t} m m}{%
  \begin{table}[#1]
    \centering
    \begin{threeparttable}
      \caption{#2}\label{#3}%
}{%
    \end{threeparttable}
  \end{table}%
}

\NewDocumentEnvironment{TableL}{O{t} m m}{%
  \begin{table*}[#1]
    \centering
    \begin{threeparttable}
      \caption{#2}\label{#3}%
}{%
    \end{threeparttable}
  \end{table*}%
}
;
% this \providecommand only guards against main.tex forgetting to,
% and matches the template's own default: the document is English.
\providecommand{\DocLanguage}{english}

% ==============================================================
% 1. PACKAGES
% ==============================================================

% --------------------------------------------------------------
% Fonts
% Times (serif, body text + math) + Helvetica (sans, headings/
% captions) + Source Code Pro (mono, verbatim/lstlisting).
% All bundled with standard TeX Live; compatible with pdfLaTeX.
\usepackage[T1]{fontenc}
\usepackage[utf8]{inputenc}
\usepackage{mathptmx}                    % Times, incl. matching math glyphs
\usepackage[scaled=0.92]{helvet}         % Helvetica
% Source Code Pro has a tall x-height and reads a size larger than
% Times beside it, so it is scaled down to match rather than left
% to dominate every \texttt and listing.
\usepackage[scaled=0.85]{sourcecodepro}

% --------------------------------------------------------------
% Math and symbols
\usepackage{amsmath}
\usepackage{amssymb}
\usepackage{amsthm}
\usepackage{textcomp}

% --------------------------------------------------------------
% Language, references and datetime
% cleveref and hyperref load last, after babel, per their
% required order: babel -> cleveref -> hyperref.
% Both languages are loaded, and \DocLanguage picks the MAIN one.
% "main=" is what actually decides the document language: babel
% otherwise takes the last language in the option list as main and
% re-asserts it at \begin{document}, overriding any preamble
% \selectlanguage. That is what kept captions reading "Figura" and
% the footer date in Portuguese under \DocLanguage=english.
% The option list has to be built by \edef first, because package
% options are not expanded when \usepackage reads them.
\edef\RBRHbabelopts{main=\DocLanguage,english,portuguese}
\expandafter\PassOptionsToPackage\expandafter{\RBRHbabelopts}{babel}
\usepackage{babel}
\usepackage[useregional]{datetime2}
\usepackage{csquotes}
\usepackage{hyperref}
\usepackage{cleveref}

% --------------------------------------------------------------
% Colors
\usepackage[dvipsnames, table]{xcolor}

% --------------------------------------------------------------
% Layout
\usepackage{geometry}
\usepackage{setspace}
\usepackage{titlesec}
\usepackage{fancyhdr}
\usepackage{lastpage}
\usepackage{paracol}                     % the two-column "sidebar + content" skeleton

% --------------------------------------------------------------
% Bibliography
% APA style with author-year citations, biber backend. Authors
% cite with \textcite{} in prose and \parencite{} in parentheses.
\usepackage[backend=biber, style=apa, citestyle=authoryear]{biblatex}

% --------------------------------------------------------------
% Sections
\usepackage{etoolbox}

% --------------------------------------------------------------
% Text
\usepackage{lipsum}                      % placeholder text only; not needed once authors write

% --------------------------------------------------------------
% Captions
\usepackage{caption}

% --------------------------------------------------------------
% Figures
\usepackage{graphicx}

% --------------------------------------------------------------
% Tables
\usepackage{array}
\usepackage{booktabs}
\usepackage{threeparttable}
% tabularx so a table can be told to occupy exactly the width it
% is given. threeparttable sizes the caption and the notes to the
% width of the tabular, so a table that fills \linewidth is what
% keeps both flush with the text column instead of ending short.
\usepackage{tabularx}
% Left-aligned flexible column: the X column, minus the centring.
\newcolumntype{L}{>{\raggedright\arraybackslash}X}

% --------------------------------------------------------------
% Boxes
\usepackage{tcolorbox}
\tcbuselibrary{skins, breakable}

% --------------------------------------------------------------
% Listings
\usepackage{listings}

% --------------------------------------------------------------
% Misc
\usepackage{fontawesome5}                % admonition icons
% ORCID mark for the author list. Draws the official logo as
% vector paths rather than embedding an SVG, which pdfLaTeX
% cannot do without --shell-escape. It scales the mark to the
% current font's X-height, so it sits inside the line and cannot
% change the leading, and it makes the icon itself the hyperlink.
\usepackage{orcidlink}

% ==============================================================
% 2. LANGUAGE SWITCH
% ==============================================================
% Nothing to do here: the document language is set by babel's
% "main=" option in section 1, which is the only thing that
% survives \begin{document}. It drives hyphenation, babel-aware
% dates (datetime2 [useregional]), quoting (csquotes) and the
% caption words ("Figure"/"Figura", "Table"/"Tabela").
%
% A preamble \selectlanguage used to sit here and did nothing:
% babel overrode it on entering the document body.

% ==============================================================
% 3. TRANSLATIONS
% ==============================================================
% Journal-facing labels, in Portuguese or English depending on
% \DocLanguage. Authors reference these by name (e.g. \trAbstract)
% rather than hardcoding either language, so the whole document
% switches with one macro.
% \ifx rather than etoolbox's \ifdefstring: \ifx is expandable, so
% \trlang survives \MakeUppercase. With \ifdefstring the conditional
% gets mangled during uppercasing and the wrong branch wins -- that
% is what made the sidebar print "ESTUDO DE CASO" under english.
%
% Both sides are normalised into plain macros first. \ifx compares
% meanings, and \newcommand makes \DocLanguage \long while a bare
% \def would not -- comparing them directly fails on that prefix
% alone and silently selects Portuguese. \edef flattens whatever
% the author wrote into the same shape as \RBRHenglish.
%
% No "@" in these names: this is a .tex file, not a .sty, so @ is
% not a letter here and \RBRHxx@yy would parse as \RBRHxx followed
% by stray characters.
\edef\RBRHlang{\DocLanguage}
\def\RBRHenglish{english}
\newcommand{\trlang}[2]{\ifx\RBRHlang\RBRHenglish #1\else #2\fi}

% --- Front matter ---
\newcommand{\trAbstract}{\trlang{Abstract}{Resumo}}
\newcommand{\trKeywords}{\trlang{Keywords}{Palavras-chave}}
\newcommand{\trListOfBoxes}{\trlang{Boxes}{Destaques}}
\newcommand{\trHighlights}{\trlang{Highlights}{Destaques do estudo}}
\newcommand{\trCorrespondence}{\trlang{Correspondence to}{Correspond\^encia para}}
\newcommand{\trCitation}{\trlang{Citation}{Como citar}}

% --- Cover sidebar ---
\newcommand{\trEditors}{\trlang{Editors}{Editores}}
\newcommand{\trEditorInChief}{\trlang{Editor-in-Chief}{Editor-chefe}}
\newcommand{\trAssociatedEditor}{\trlang{Associated Editor}{Editor associado}}
\newcommand{\trCompetingInterests}{\trlang{Competing interests}{Conflito de interesses}}
\newcommand{\trDataAvailability}{\trlang{Data availability}{Disponibilidade dos dados}}
\newcommand{\trPublicationDates}{\trlang{Publication dates}{Datas de publica\c{c}\~ao}}
\newcommand{\trPublished}{\trlang{Published}{Publicado em}}
\newcommand{\trIssue}{\trlang{Issue}{Edi\c{c}\~ao}}
\newcommand{\trHistory}{\trlang{History}{Hist\'orico}}
\newcommand{\trReceived}{\trlang{Received}{Recebido em}}
\newcommand{\trReviewed}{\trlang{Reviewed}{Revisado em}}
\newcommand{\trAccepted}{\trlang{Accepted}{Aceito em}}

% --- Back matter ---
\newcommand{\trReferences}{\trlang{References}{Refer\^encias}}
\newcommand{\trAppendix}{\trlang{Appendix}{Ap\^endice}}
\newcommand{\trAppendices}{\trlang{Appendices}{Ap\^endices}}
\newcommand{\trAcknowledgements}{\trlang{Acknowledgements}{Agradecimentos}}

% --- Items ---
\newcommand{\trListing}{\trlang{Code}{C\'odigo}}
\newcommand{\trListings}{\trlang{Code}{C\'odigos}}
\newcommand{\trDraft}{\trlang{DRAFT}{MINUTA}}

% --- Admonitions ---
\newcommand{\trAdmNote}{\trlang{Note}{Nota}}
\newcommand{\trAdmTip}{\trlang{Tip}{Dica}}
\newcommand{\trAdmImportant}{\trlang{Important}{Importante}}
\newcommand{\trAdmWarning}{\trlang{Warning}{Alerta}}
\newcommand{\trAdmCaution}{\trlang{Caution}{Cuidado}}

% ==============================================================
% 4. COLORS
% ==============================================================
% Override this section to apply a different brand/theme.

% --- Grayscale utilities ---
\definecolor{verylightgray}{gray}{0.97}
\definecolor{lightgray}{gray}{0.85}
\definecolor{midgray}{gray}{0.50}
\definecolor{heavygray}{gray}{0.3}

% --- Brand / accent ---
% The journal's blue is the accent colour: it carries the article
% title, the section headings, the abstract headings and the
% metadata column. The masthead itself is set in black, so the
% blue marks structure rather than branding.
%
% Everything already keyed to AccentColor follows from here --
% caption labels, theorem heads and the tcolorbox palette, whose
% tints derive from it.
\definecolor{AccentColor}{RGB}{38,52,139}
\colorlet{AccentColorVeryLight}{AccentColor!5}
\colorlet{AccentColorLight}{AccentColor!20}

% --- Semantic aliases ---
\definecolor{HeadingColor}{named}{AccentColor}
\definecolor{LinkColor}{RGB}{38,52,139}
\definecolor{RuleColor}{named}{midgray}
\definecolor{TextColor}{named}{black}

% --- Page and text ---
\pagecolor{white}
\color{TextColor}

% --- Links ---
\hypersetup{
    colorlinks=true,
    linkcolor=LinkColor,
    citecolor=LinkColor,
    filecolor=LinkColor,
    urlcolor=LinkColor,
}

% ==============================================================
% 5. HEADER AND FOOTER
% ==============================================================
% Every field comes from cover.tex, which build_cover.py generates
% from the TOML files. Both journal logos sit in the running
% header at the same height, one at each margin; the footer
% carries the short citation on the left and the page count on
% the right.
\newcommand{\RBRHlogoheight}{8mm}
\newcommand{\HeaderLeft}{\includegraphics[height=\RBRHlogoheight]{\JournalLogoLeft}}
\newcommand{\HeaderRight}{\includegraphics[height=\RBRHlogoheight]{\JournalLogoRight}}
\newcommand{\FooterLeft}{\ArticleShortCitation}
\newcommand{\FooterCenter}{}
\newcommand{\FooterRight}{\thepage\;/\;\pageref{LastPage}}

\fancyhfoffset[L]{0mm}
\fancyhfoffset[R]{0mm}

\fancypagestyle{plain}{%
  \fancyhf{}
  \fancyhead[L]{\HeaderLeft}
  \fancyhead[R]{\HeaderRight}
  \renewcommand{\headrulewidth}{0.4pt}
  \fancyfoot[L]{\sffamily\scriptsize{\FooterLeft}}
  \fancyfoot[C]{\sffamily\scriptsize{\FooterCenter}}
  \fancyfoot[R]{\sffamily\scriptsize{\FooterRight}}
  \renewcommand{\footrulewidth}{0.2pt}
}

\pagestyle{fancy}
\fancyhf{}
\fancyhead[L]{\HeaderLeft}
\fancyhead[R]{\HeaderRight}
\renewcommand{\headrulewidth}{0.4pt}
\fancyfoot[L]{\sffamily\scriptsize{\FooterLeft}}
\fancyfoot[C]{\sffamily\scriptsize{\FooterCenter}}
\fancyfoot[R]{\sffamily\scriptsize{\FooterRight}}
\renewcommand{\footrulewidth}{0.2pt}

% ==============================================================
% 6. PAGE LAYOUT AND SPACING
% ==============================================================
\geometry{
 a4paper,
 left=20mm,
 right=20mm,
 top=20mm,
 includefoot,
 bottom=20mm,
 headheight=15mm,
 headsep=5mm,
 footskip=10mm,
}

\newcommand{\pagesectionvspace}{\vspace*{0mm}}

\setlength{\parindent}{0mm}
\setlength{\parskip}{2mm}

\setlength{\intextsep}{1mm}
\setlength{\textfloatsep}{2mm}
\setlength{\floatsep}{2mm}
\setlength{\abovecaptionskip}{1mm}
\setlength{\belowcaptionskip}{0mm}

\linespread{1}

\titlespacing*{\section}      {0mm}{4mm}{0mm}
\titlespacing*{\subsection}   {0mm}{2mm}{0mm}
\titlespacing*{\subsubsection}{0mm}{2mm}{0mm}
\titlespacing{\paragraph}{0mm}{2mm}{2mm}

% --------------------------------------------------------------
% Camera-ready sidebar column (left, narrow) + content column
% (right, wide). Width tuned for the DOI/highlights/correspondence
% sidebar used by main.tex; adjust \RBRHColumnRatio to change it.
\newcommand{\RBRHColumnRatio}{0.258}

% Lengths for the small-with-side-caption figure pattern (see
% contents.tex): a small, fixed-size image whose caption spans
% into the space the narrow sidebar column would otherwise
% occupy. Allocated once here; reused via \setlength (not
% \newlength) at the use site.
\newlength{\sidecapimgwidth}
\newlength{\sidecapgap}
\newlength{\sidecapwidth}
\newlength{\smallfigfixedwidth}

% ==============================================================
% 7. BIBLIOGRAPHY RESOURCES
% ==============================================================
\addbibresource{references.bib}
% Reference list one size down from the body text.
\renewcommand*{\bibfont}{\small}

% ==============================================================
% 8. CAPTIONS
% ==============================================================
\DeclareCaptionLabelSeparator{pipe}{\;|\;}

\captionsetup[figure]{
    font={footnotesize, sf, color=heavygray},
    labelfont={bf, color=AccentColor},
    labelsep=quad,
}

% No "width=" here. Inside paracol \textwidth is still the full
% page width, not the width of the column being typeset, so
% width=\textwidth made table captions overrun the text column and
% spill into the sidebar. Left unset, the caption takes the width
% of whatever contains it -- the text column for TableM, the full
% span for TableL -- which is right in both cases.
\captionsetup[table]{
    font={footnotesize, sf},
    labelfont={bf, color=AccentColor},
    labelsep=pipe,
}

\DeclareCaptionStyle{lstlisting}{
    font={footnotesize, sf},
    labelfont={bf, color=AccentColor},
    labelsep=pipe,
    justification=centering
}
\captionsetup[lstlisting]{style=lstlisting}

% All captions: justified, flush left, regardless of length.
% singlelinecheck=false disables caption's default of auto-centering
% any caption that happens to fit on a single line.
\captionsetup{justification=justified,singlelinecheck=false,labelfont=bf}

% ==============================================================
% 9. SECTION TITLES
% ==============================================================
\titleformat{\section}
  {\bfseries\sffamily\Large\color{HeadingColor}}
  {\makebox[8mm][l]{\thesection}}{2mm}{}

\titleformat{\subsection}
  {\bfseries\sffamily\large\color{HeadingColor}}
  {\makebox[8mm][l]{\thesubsection}}{2mm}{}

\titleformat{\subsubsection}
  {\sffamily\normalsize\color{HeadingColor}}
  {\makebox[10mm][l]{\thesubsubsection}}{4mm}{}

\titleformat{\paragraph}[runin]
  {\bfseries\sffamily\normalsize\color{HeadingColor}}
  {}{0mm}{}[.]

% Shared heading for auxiliary lists (References, List of Boxes,
% Glossary, ...) so they match the plain \section* styling used
% by \tableofcontents / \listoffigures / \listoftables, instead of
% falling back to each package's own internal heading.
\newcommand{\auxlisttitle}[1]{%
  \begingroup
    \titlespacing*{\section}{0mm}{6mm}{2mm}%
    \section*{#1}%
  \endgroup
  \phantomsection
  \addcontentsline{toc}{section}{#1}%
}

% ==============================================================
% 10. THEOREM-LIKE ENVIRONMENTS
% ==============================================================
\providecommand{\theoremname}{\trlang{Theorem}{Teorema}}
\providecommand{\lemmaname}{\trlang{Lemma}{Lema}}
\providecommand{\corollaryname}{\trlang{Corollary}{Corol\'ario}}
\providecommand{\definitionname}{\trlang{Definition}{Defini\c{c}\~ao}}
\providecommand{\examplename}{\trlang{Example}{Exemplo}}
\providecommand{\remarkname}{\trlang{Remark}{Observa\c{c}\~ao}}

\makeatletter
\renewcommand{\th@plain}{%
    \thm@headfont{\bfseries\sffamily\color{AccentColor}}%
    \itshape
}
\renewcommand{\th@definition}{%
    \thm@headfont{\bfseries\sffamily\color{AccentColor}}%
    \normalfont
}
\renewcommand{\th@remark}{%
    \thm@headfont{\itshape\sffamily\color{AccentColor}}%
    \normalfont
}
\makeatother

\newtheorem{theorem}{\theoremname}
\newtheorem{lemma}[theorem]{\lemmaname}
\newtheorem{corollary}[theorem]{\corollaryname}
\newtheorem{definition}{\definitionname}
\newtheorem{example}{\examplename}
\newtheorem*{remark}{\remarkname}

% ==============================================================
% 11. LISTINGS
% ==============================================================
\renewcommand{\lstlistingname}{\trListing}
\crefname{lstlisting}{\trListing}{\trListings}
\Crefname{lstlisting}{\trListing}{\trListings}

\lstdefinestyle{codeblock}{
    backgroundcolor=\color{black!5},
    commentstyle=\color{black!40}\itshape,
    keywordstyle=\color{purple}\bfseries,
    stringstyle=\color{teal},
    numberstyle=\tiny\color{black!40},
    basicstyle=\ttfamily\footnotesize,
    breaklines=true,
    breakatwhitespace=true,
    postbreak=\mbox{\textcolor{gray}{$\hookrightarrow$}\space},
    numbersep=8pt,
    tabsize=4,
    keepspaces=true,
    frame=none,
    captionpos=b,
    showspaces=false,
    showstringspaces=false,
    showtabs=false,
    numbers=left,
}
\lstset{style=codeblock, language=Python}
\lstset{basicstyle=\ttfamily\footnotesize}
\renewcommand{\ttdefault}{SourceCodePro-TLF}
\renewcommand{\familydefault}{\rmdefault}

% ==============================================================
% 12. BOXES
% ==============================================================
\tcbset{
    boxbasestyle/.style={
        enhanced,
        breakable,
        colback      = AccentColorVeryLight,
        colbacktitle = AccentColorLight,
        coltitle     = AccentColor,
        colframe     = AccentColor,
        boxrule      = 0.5pt,
        arc          = 2pt,
        fonttitle    = \bfseries\sffamily\small,
        fontupper    = \sffamily\footnotesize,
        boxsep       = 0mm,
        top          = 2mm,
        bottom       = 2mm,
        left         = 2mm,
        right        = 2mm,
        toptitle     = 1mm,
        bottomtitle  = 1mm,
    }
}

% boxstory — numbered, breakable, document-wide counter.
% Usage: \begin{boxstory}{Title} ... \end{boxstory}
\newtcolorbox[
    auto counter,
    list inside = lob,
]{boxstory}[2][]{
    boxbasestyle,
    title = {\sffamily Box~\thetcbcounter\quad #2},
    #1,
}

\newcommand{\notitleheading}[1]{}
\newcommand{\listofboxes}{%
    \auxlisttitle{\trListOfBoxes}%
    \tcblistof[\notitleheading]{lob}{\trListOfBoxes}%
}

% ==============================================================
% 13. ADMONITIONS
% ==============================================================
% Color-coded unnumbered admonition boxes. Icons require fontawesome5.
\newtcolorbox{amdnote}{
    boxbasestyle,
    colback = cyan!5, colbacktitle = cyan!50!black, colframe = cyan!50!black,
    coltitle = white, title = {\faInfoCircle\quad \trAdmNote},
}
\newtcolorbox{amdtip}{
    boxbasestyle,
    colback = green!5!white, colbacktitle = green!40!black, colframe = green!40!black,
    coltitle = white, title = {\faLightbulb\quad \trAdmTip},
}
\newtcolorbox{amdimportant}{
    boxbasestyle,
    colback = violet!5, colbacktitle = violet!60!black, colframe = violet!60!black,
    coltitle = white, title = {\faExclamationCircle\quad \trAdmImportant},
}
\newtcolorbox{amdwarning}{
    boxbasestyle,
    colback = orange!5, colbacktitle = orange!70!black, colframe = orange!70!black,
    coltitle = white, title = {\faExclamationTriangle\quad \trAdmWarning},
}
\newtcolorbox{amdcaution}{
    boxbasestyle,
    colback = red!5, colbacktitle = red!70!black, colframe = red!70!black,
    coltitle = white, title = {\faBan\quad \trAdmCaution},
}

% ==============================================================
% 14. DRAFT WATERMARK
% ==============================================================
% main.tex sets \IsDraft to "true" to stamp every page "DRAFT" /
% "MINUTA" while the manuscript is a work in progress.
\providecommand{\IsDraft}{false}
\ifdefstring{\IsDraft}{true}{
    \usepackage{draftwatermark}
    \SetWatermarkText{\trDraft}
    \SetWatermarkScale{1}
    \SetWatermarkColor[gray]{0.8}
    \SetWatermarkAngle{30}
}{}

% ==============================================================
% 15. COVER
% ==============================================================
% Everything here reads macros defined in cover.tex, which
% build_cover.py generates from journal.toml, article.toml and
% authors.toml. Nothing on the cover is written by hand, which is
% why contents.tex can begin straight at the first section.
%
% \PrintMasthead fills the wide column: journal name, title,
% authors, affiliations, then both abstracts. \PrintSidebar fills
% the narrow one, in the order RBRH prints it.

% The author and affiliation blocks are long runs of proper nouns
% and institution names, which offer TeX very few legal break
% points. On the default settings the author list overflowed the
% column by about 10pt.
%
% Justified was tried first and does fix the overflow, with
% \tolerance raised and \emergencystretch added -- but it pays for
% it by hyphenating personal names across lines ("An-derson",
% "Fer-nando") and by leaving one very loose line. Breaking a
% surname is not acceptable in a masthead, so these blocks are set
% ragged right instead, with hyphenation switched off outright.
% The right edge is uneven; nothing overflows and no name breaks.
%
% To go back to justified, swap \raggedright and the two penalties
% for: \tolerance=4000 \emergencystretch=2em
%
% These are paragraph-breaking parameters, so they only take
% effect if the paragraph ends while they are still in scope --
% hence the \par inside each group below.
\newcommand{\RBRHnameblock}{%
  \raggedright
  \hyphenpenalty=10000
  \exhyphenpenalty=10000
}

% --- Wide column ---
\newcommand{\PrintMasthead}{%
  {\noindent \sffamily \Large\textbf{\JournalName}}\\
  {\noindent \large \textit{\JournalNameEn}}\\[3mm]
  \begin{center}
    \raggedright {\Large \sffamily \textbf{\color{AccentColor}\ArticleTitle}}\\[3mm]
  \end{center}
  {\RBRHnameblock\noindent\AuthorBlock\par}
  \vspace{2mm}
  % Affiliations in the same grey as the captions, so they read as
  % secondary to the author names without fading out like the
  % licence notice does.
  {\footnotesize\color{heavygray}\RBRHnameblock\noindent\AffiliationBlock\par}
  \vspace{2mm}
  \hrule
  \vspace{2mm}
  \PrintAbstracts
}

% Both languages are printed whatever \DocLanguage says: RBRH
% carries an English abstract and a Portuguese resumo regardless
% of the language the body is written in.
% The \small applies to the heading words only -- "Abstract",
% "Keywords:" and their Portuguese counterparts, which were
% sitting at full body size. The abstract text itself stays at
% \footnotesize.
\newcommand{\PrintAbstracts}{%
  \begin{spacing}{0.8}
  \small
  {\noindent \sffamily\textbf{\color{AccentColor}Abstract}\quad{\footnotesize\color{TextColor} \AbstractEn}}\\[1mm]
  {\noindent \sffamily\textbf{\color{AccentColor}Keywords:}\quad{\footnotesize\color{TextColor} \KeywordsEn}}\\[3mm]
  {\noindent \sffamily\textbf{\color{AccentColor}Resumo}\quad{\footnotesize\color{TextColor} \AbstractPt}}\\[1mm]
  {\noindent \sffamily\textbf{\color{AccentColor}Palavras-chave:}\quad{\footnotesize\color{TextColor} \KeywordsPt}}
  \end{spacing}
}

% --- Narrow column ---
% One helper per block, so the running order is readable and can
% be reshuffled by editing \PrintSidebar alone.
% The accent marks the structure of the column -- the heading of
% each block -- while the content under it stays in body colour.
\newcommand{\RBRHcoverblock}[2]{%
  {\color{AccentColor}\textbf{#1}}\\[1mm]
  #2
  \\[4mm]
}

\newcommand{\PrintSidebar}{%
  {
  \sffamily
  \scriptsize
  \raggedright
  % licence block at the foot overrides it with grey.
  \textbf{\MakeUppercase{\ArticleType}}\\
  \ifdefempty{\ArticleDOI}{}{DOI\;\ArticleDOI\\[4mm]}
  \RBRHcoverblock{\trHighlights:}{\ArticleHighlights}
  \RBRHcoverblock{\trCitation:}{\ArticleCitation}
  \RBRHcoverblock{\trCorrespondence:}{\CorrespondingBlock}
  \RBRHcoverblock{\trCompetingInterests:}{\CompetingInterests}
  \RBRHcoverblock{\trAcknowledgements:}{\Acknowledgements}
  \RBRHcoverblock{\trEditors:}{%
    \trEditorInChief:\\ \EditorInChief\\[1mm]
    \trAssociatedEditor:\\ \AssociatedEditor}
  \RBRHcoverblock{\trDataAvailability:}{\DataAvailability}
  \RBRHcoverblock{\trPublicationDates:}{%
    \trPublished\ \DatePublished\\
    \trIssue\ \ArticleIssue}
  \RBRHcoverblock{\trHistory:}{%
    \trReceived\ \DateReceived\\
    \trReviewed\ \DateReviewed\\
    \trAccepted\ \DateAccepted}
  % Push the licence to the foot of the column: everything above
  % sets at its natural height and this absorbs the slack, so the
  % notice sits on the baseline of the page however much or little
  % metadata precedes it.
  \vfill
  \PrintLicense
  }
}

% CC-BY notice, set in grey so it reads as a footer to the
% sidebar rather than as content. The badge is a vector PDF
% converted from assets/cc-by.svg.
\newcommand{\PrintLicense}{%
  {\color{midgray}
   \includegraphics[height=3.2mm]{\LicenseBadge}\\[1mm]
   \LicenseCopyright\\
   \LicenseStatement\par}
}

% ==============================================================
% 16. APPENDIX RENUMBERING
% ==============================================================
% Call \BeginAppendix in contents.tex to switch into lettered,
% restarted numbering (A, A.1, Figure A.1, ...) for supplementary
% material.
\newcommand{\BeginAppendix}{%
  \appendix
  \renewcommand{\thesection}{\Alph{section}}
  \renewcommand{\thesubsection}{\Alph{section}.\arabic{subsection}}
  \renewcommand{\thesubsubsection}{\Alph{section}.\arabic{subsection}.\arabic{subsubsection}}
  \setcounter{section}{0}
  \setcounter{equation}{0}
  \setcounter{figure}{0}
  \setcounter{table}{0}
  \setcounter{lstlisting}{0}
  \renewcommand{\theequation}{\Alph{section}.\arabic{equation}}
  \renewcommand{\thefigure}{\Alph{section}.\arabic{figure}}
  \renewcommand{\thetable}{\Alph{section}.\arabic{table}}
  \renewcommand{\thelstlisting}{\Alph{section}.\arabic{lstlisting}}
  \pagesectionvspace
  \auxlisttitle{\trAppendices}
  \label{sec:appendix}
  \vspace{10mm}
  \small
  % Headings one level down too. The body is already \small, which
  % left \Large section titles looking oversized against it.
  \titleformat{\section}
    {\bfseries\sffamily\large\color{HeadingColor}}
    {\makebox[8mm][l]{\thesection}}{2mm}{}
  \titleformat{\subsection}
    {\bfseries\sffamily\normalsize\color{HeadingColor}}
    {\makebox[8mm][l]{\thesubsection}}{2mm}{}
  \titleformat{\subsubsection}
    {\sffamily\small\color{HeadingColor}}
    {\makebox[10mm][l]{\thesubsubsection}}{4mm}{}
}

% ==============================================================
% 17. SHARED COUNTERS ACROSS PARACOL COLUMNS
% ==============================================================
% paracol makes every counter column-LOCAL by default -- its
% internal list of global counters ships as just {page}. Each
% column therefore keeps its own `figure`, `table` and `equation`.
% A figure written inside a column and one written in a
% \switchcolumn[1]*[...] spanning block are stepped in different
% counter states, so the spanning one restarts at 1 and the
% numbering appears to repeat. `figure` is shared instead: the two
% columns here are a metadata sidebar plus one article body, not
% parallel independent texts.
%
% Name the counter explicitly -- do NOT use \globalcounter*, and
% do NOT add `table` here. Both were tried and both break this
% document with "Missing \endcsname inserted" at the first
% \section (the starred form because it sweeps in every counter,
% including `table`). Bisected against this exact package set:
%   figure   -> safe, and is what fixes the numbering bug
%   equation -> safe, but unused; add only if an equation ever
%               goes inside a \switchcolumn[1]*[...] span
%   table    -> BREAKS the build, do not enable
% `footnote` is left alone too: paracol manages its scope itself.
\globalcounter{figure}

% ==============================================================
% 18. AUTHOR-FACING FIGURE AND TABLE MACROS
% ==============================================================
% So contents.tex never carries layout plumbing:
%
%   \FigureM[pos]{name}{caption}{label}        text width
%   \FigureS[pos][frac]{name}{caption}{label}  small -- frac of the
%                                              text width, default
%                                              0.67 -- and the only
%                                              size whose caption
%                                              sits beside the
%                                              image, spanning into
%                                              the sidebar column
%   \FigureL[pos]{name}{caption}{label}        full width, margin to
%                                              margin, matching the
%                                              width of the header
%                                              rule
%
%   \begin{TableM}[pos]{caption}{label} ... \end{TableM}  text width
%   \begin{TableL}[pos]{caption}{label} ... \end{TableL}  full width
%
% [pos] is the usual LaTeX float placement specifier and defaults
% to t (top). The table environments take the tabular material as
% their body and put the caption above it, per table convention;
% figures caption below.
%
% <name> is the bare file name: no folder, no extension. The file
% is resolved inside figures/, preferring .pdf, then .png, then
% .jpeg, then .jpg. A name that matches nothing degrades to a
% placeholder box rather than failing the build.
%
% M/S/L map onto real floats -- figure/table for M, and the
% double-column figure*/table* for S and L, which is what makes
% them span the sidebar. paracol supports both (it carries its own
% \pcol@floatplacement and spanning-float machinery), and the
% figure counter is shared across columns by section 19 above, so
% numbering follows source order. \label goes straight after
% \caption inside the float, where \@currentlabel is correct.

\newif\ifRBRHfigfound
\newcommand{\RBRHfigpath}{}

% Resolve figures/<name> against the supported extensions, in
% order of preference.
\newcommand{\RBRHfindfigure}[1]{%
  \IfFileExists{figures/#1.pdf}%
    {\RBRHfigfoundtrue \def\RBRHfigpath{figures/#1.pdf}}%
    {\IfFileExists{figures/#1.png}%
      {\RBRHfigfoundtrue \def\RBRHfigpath{figures/#1.png}}%
      {\IfFileExists{figures/#1.jpeg}%
        {\RBRHfigfoundtrue \def\RBRHfigpath{figures/#1.jpeg}}%
        {\IfFileExists{figures/#1.jpg}%
          {\RBRHfigfoundtrue \def\RBRHfigpath{figures/#1.jpg}}%
          {\RBRHfigfoundfalse \def\RBRHfigpath{}}}}}%
}

% Placeholder for a missing file: gray box in golden ratio
% (height = width / 1.618) carrying a red notice. \detokenize
% keeps names containing "_" printable.
\newlength{\RBRHfigwidth}
\newlength{\RBRHfigheight}
\newcommand{\RBRHmissingfigure}[2]{%
  \setlength{\RBRHfigwidth}{#1}%
  \setlength{\RBRHfigheight}{\dimexpr\RBRHfigwidth*1000/1618\relax}%
  \begingroup
    \setlength{\fboxsep}{0pt}%
    \colorbox{lightgray}{%
      \begin{minipage}[c][\RBRHfigheight][c]{\RBRHfigwidth}%
        \centering
        \textcolor{red}{\sffamily\small
          Figure \texttt{\detokenize{#2}} not found}%
      \end{minipage}%
    }%
  \endgroup
}

% The image at #1 wide, or the placeholder if it is missing.
\newcommand{\RBRHfigbody}[2]{%
  \ifRBRHfigfound
    \includegraphics[width=#1]{\RBRHfigpath}%
  \else
    \RBRHmissingfigure{#1}{#2}%
  \fi
}

% --- M: text width ---
\NewDocumentCommand{\FigureM}{O{t} m m m}{%
  \begin{figure}[#1]
    \centering
    \RBRHfindfigure{#2}%
    \RBRHfigbody{\linewidth}{#2}%
    \caption{#3}\label{#4}%
  \end{figure}%
}

% --- L: full width, margin to margin ---
\NewDocumentCommand{\FigureL}{O{t} m m m}{%
  \begin{figure*}[#1]
    \centering
    \RBRHfindfigure{#2}%
    \RBRHfigbody{\textwidth}{#2}%
    \caption{#3}\label{#4}%
  \end{figure*}%
}

% --- S: small, caption beside the image, spanning the sidebar ---
% The image width is captured BEFORE the float opens, while
% \linewidth still means the text column: inside a double-column
% float LaTeX widens it to \textwidth. The float body is typeset
% here at the call site, so the captured value is the right one.
\NewDocumentCommand{\FigureS}{O{t} O{0.67} m m m}{%
  \setlength{\smallfigfixedwidth}{#2\linewidth}%
  \begin{figure*}[#1]
    \RBRHfindfigure{#3}%
    \setlength{\sidecapimgwidth}{\smallfigfixedwidth}%
    \setlength{\sidecapgap}{6mm}%
    \setlength{\sidecapwidth}{\textwidth}%
    \addtolength{\sidecapwidth}{-\sidecapimgwidth}%
    \addtolength{\sidecapwidth}{-\sidecapgap}%
    \noindent
    % \vspace{0pt} as the first thing in each minipage is what makes
    % [t] align the TOPS. Without it, [t] aligns first baselines, and
    % a minipage holding just a graphic has its baseline at the
    % graphic's bottom -- which dropped the caption to the foot of
    % the image instead of level with its top. \abovecaptionskip is
    % zeroed for the same reason: it is meant to separate a figure
    % from a caption printed below it, not beside it.
    \begin{minipage}[t]{\sidecapwidth}
      \vspace{0pt}%
      \setlength{\abovecaptionskip}{0pt}%
      \caption{#4}\label{#5}%
    \end{minipage}%
    \hspace{\sidecapgap}%
    \begin{minipage}[t]{\sidecapimgwidth}
      \vspace{0pt}%
      \RBRHfigbody{\linewidth}{#3}%
    \end{minipage}%
  \end{figure*}%
}

% --- Tables: caption above the body, per table convention ---
% The body is wrapped in threeparttable, so the caption and any
% notes take the width of the tabular rather than running wider,
% and authors can add a notes block without extra scaffolding:
%
%   \begin{TableM}{caption}{label}
%     \begin{tabular}{ll} ... \end{tabular}
%     \begin{tablenotes}
%       \item Note text.
%     \end{tablenotes}
%   \end{TableM}
\NewDocumentEnvironment{TableM}{O{t} m m}{%
  \begin{table}[#1]
    \centering
    \begin{threeparttable}
      \caption{#2}\label{#3}%
}{%
    \end{threeparttable}
  \end{table}%
}

\NewDocumentEnvironment{TableL}{O{t} m m}{%
  \begin{table*}[#1]
    \centering
    \begin{threeparttable}
      \caption{#2}\label{#3}%
}{%
    \end{threeparttable}
  \end{table*}%
}
;
% this \providecommand only guards against main.tex forgetting to,
% and matches the template's own default: the document is English.
\providecommand{\DocLanguage}{english}

% ==============================================================
% 1. PACKAGES
% ==============================================================

% --------------------------------------------------------------
% Fonts
% Times (serif, body text + math) + Helvetica (sans, headings/
% captions) + Source Code Pro (mono, verbatim/lstlisting).
% All bundled with standard TeX Live; compatible with pdfLaTeX.
\usepackage[T1]{fontenc}
\usepackage[utf8]{inputenc}
\usepackage{mathptmx}                    % Times, incl. matching math glyphs
\usepackage[scaled=0.92]{helvet}         % Helvetica
% Source Code Pro has a tall x-height and reads a size larger than
% Times beside it, so it is scaled down to match rather than left
% to dominate every \texttt and listing.
\usepackage[scaled=0.85]{sourcecodepro}

% --------------------------------------------------------------
% Math and symbols
\usepackage{amsmath}
\usepackage{amssymb}
\usepackage{amsthm}
\usepackage{textcomp}

% --------------------------------------------------------------
% Language, references and datetime
% cleveref and hyperref load last, after babel, per their
% required order: babel -> cleveref -> hyperref.
% Both languages are loaded, and \DocLanguage picks the MAIN one.
% "main=" is what actually decides the document language: babel
% otherwise takes the last language in the option list as main and
% re-asserts it at \begin{document}, overriding any preamble
% \selectlanguage. That is what kept captions reading "Figura" and
% the footer date in Portuguese under \DocLanguage=english.
% The option list has to be built by \edef first, because package
% options are not expanded when \usepackage reads them.
\edef\RBRHbabelopts{main=\DocLanguage,english,portuguese}
\expandafter\PassOptionsToPackage\expandafter{\RBRHbabelopts}{babel}
\usepackage{babel}
\usepackage[useregional]{datetime2}
\usepackage{csquotes}
\usepackage{hyperref}
\usepackage{cleveref}

% --------------------------------------------------------------
% Colors
\usepackage[dvipsnames, table]{xcolor}

% --------------------------------------------------------------
% Layout
\usepackage{geometry}
\usepackage{setspace}
\usepackage{titlesec}
\usepackage{fancyhdr}
\usepackage{lastpage}
\usepackage{paracol}                     % the two-column "sidebar + content" skeleton

% --------------------------------------------------------------
% Bibliography
% APA style with author-year citations, biber backend. Authors
% cite with \textcite{} in prose and \parencite{} in parentheses.
\usepackage[backend=biber, style=apa, citestyle=authoryear]{biblatex}

% --------------------------------------------------------------
% Sections
\usepackage{etoolbox}

% --------------------------------------------------------------
% Text
\usepackage{lipsum}                      % placeholder text only; not needed once authors write

% --------------------------------------------------------------
% Captions
\usepackage{caption}

% --------------------------------------------------------------
% Figures
\usepackage{graphicx}

% --------------------------------------------------------------
% Tables
\usepackage{array}
\usepackage{booktabs}
\usepackage{threeparttable}
% tabularx so a table can be told to occupy exactly the width it
% is given. threeparttable sizes the caption and the notes to the
% width of the tabular, so a table that fills \linewidth is what
% keeps both flush with the text column instead of ending short.
\usepackage{tabularx}
% Left-aligned flexible column: the X column, minus the centring.
\newcolumntype{L}{>{\raggedright\arraybackslash}X}

% --------------------------------------------------------------
% Boxes
\usepackage{tcolorbox}
\tcbuselibrary{skins, breakable}

% --------------------------------------------------------------
% Listings
\usepackage{listings}

% --------------------------------------------------------------
% Misc
\usepackage{fontawesome5}                % admonition icons
% ORCID mark for the author list. Draws the official logo as
% vector paths rather than embedding an SVG, which pdfLaTeX
% cannot do without --shell-escape. It scales the mark to the
% current font's X-height, so it sits inside the line and cannot
% change the leading, and it makes the icon itself the hyperlink.
\usepackage{orcidlink}

% ==============================================================
% 2. LANGUAGE SWITCH
% ==============================================================
% Nothing to do here: the document language is set by babel's
% "main=" option in section 1, which is the only thing that
% survives \begin{document}. It drives hyphenation, babel-aware
% dates (datetime2 [useregional]), quoting (csquotes) and the
% caption words ("Figure"/"Figura", "Table"/"Tabela").
%
% A preamble \selectlanguage used to sit here and did nothing:
% babel overrode it on entering the document body.

% ==============================================================
% 3. TRANSLATIONS
% ==============================================================
% Journal-facing labels, in Portuguese or English depending on
% \DocLanguage. Authors reference these by name (e.g. \trAbstract)
% rather than hardcoding either language, so the whole document
% switches with one macro.
% \ifx rather than etoolbox's \ifdefstring: \ifx is expandable, so
% \trlang survives \MakeUppercase. With \ifdefstring the conditional
% gets mangled during uppercasing and the wrong branch wins -- that
% is what made the sidebar print "ESTUDO DE CASO" under english.
%
% Both sides are normalised into plain macros first. \ifx compares
% meanings, and \newcommand makes \DocLanguage \long while a bare
% \def would not -- comparing them directly fails on that prefix
% alone and silently selects Portuguese. \edef flattens whatever
% the author wrote into the same shape as \RBRHenglish.
%
% No "@" in these names: this is a .tex file, not a .sty, so @ is
% not a letter here and \RBRHxx@yy would parse as \RBRHxx followed
% by stray characters.
\edef\RBRHlang{\DocLanguage}
\def\RBRHenglish{english}
\newcommand{\trlang}[2]{\ifx\RBRHlang\RBRHenglish #1\else #2\fi}

% --- Front matter ---
\newcommand{\trAbstract}{\trlang{Abstract}{Resumo}}
\newcommand{\trKeywords}{\trlang{Keywords}{Palavras-chave}}
\newcommand{\trListOfBoxes}{\trlang{Boxes}{Destaques}}
\newcommand{\trHighlights}{\trlang{Highlights}{Destaques do estudo}}
\newcommand{\trCorrespondence}{\trlang{Correspondence to}{Correspond\^encia para}}
\newcommand{\trCitation}{\trlang{Citation}{Como citar}}

% --- Cover sidebar ---
\newcommand{\trEditors}{\trlang{Editors}{Editores}}
\newcommand{\trEditorInChief}{\trlang{Editor-in-Chief}{Editor-chefe}}
\newcommand{\trAssociatedEditor}{\trlang{Associated Editor}{Editor associado}}
\newcommand{\trCompetingInterests}{\trlang{Competing interests}{Conflito de interesses}}
\newcommand{\trDataAvailability}{\trlang{Data availability}{Disponibilidade dos dados}}
\newcommand{\trPublicationDates}{\trlang{Publication dates}{Datas de publica\c{c}\~ao}}
\newcommand{\trPublished}{\trlang{Published}{Publicado em}}
\newcommand{\trIssue}{\trlang{Issue}{Edi\c{c}\~ao}}
\newcommand{\trHistory}{\trlang{History}{Hist\'orico}}
\newcommand{\trReceived}{\trlang{Received}{Recebido em}}
\newcommand{\trReviewed}{\trlang{Reviewed}{Revisado em}}
\newcommand{\trAccepted}{\trlang{Accepted}{Aceito em}}

% --- Back matter ---
\newcommand{\trReferences}{\trlang{References}{Refer\^encias}}
\newcommand{\trAppendix}{\trlang{Appendix}{Ap\^endice}}
\newcommand{\trAppendices}{\trlang{Appendices}{Ap\^endices}}
\newcommand{\trAcknowledgements}{\trlang{Acknowledgements}{Agradecimentos}}

% --- Items ---
\newcommand{\trListing}{\trlang{Code}{C\'odigo}}
\newcommand{\trListings}{\trlang{Code}{C\'odigos}}
\newcommand{\trDraft}{\trlang{DRAFT}{MINUTA}}

% --- Admonitions ---
\newcommand{\trAdmNote}{\trlang{Note}{Nota}}
\newcommand{\trAdmTip}{\trlang{Tip}{Dica}}
\newcommand{\trAdmImportant}{\trlang{Important}{Importante}}
\newcommand{\trAdmWarning}{\trlang{Warning}{Alerta}}
\newcommand{\trAdmCaution}{\trlang{Caution}{Cuidado}}

% ==============================================================
% 4. COLORS
% ==============================================================
% Override this section to apply a different brand/theme.

% --- Grayscale utilities ---
\definecolor{verylightgray}{gray}{0.97}
\definecolor{lightgray}{gray}{0.85}
\definecolor{midgray}{gray}{0.50}
\definecolor{heavygray}{gray}{0.3}

% --- Brand / accent ---
% The journal's blue is the accent colour: it carries the article
% title, the section headings, the abstract headings and the
% metadata column. The masthead itself is set in black, so the
% blue marks structure rather than branding.
%
% Everything already keyed to AccentColor follows from here --
% caption labels, theorem heads and the tcolorbox palette, whose
% tints derive from it.
\definecolor{AccentColor}{RGB}{38,52,139}
\colorlet{AccentColorVeryLight}{AccentColor!5}
\colorlet{AccentColorLight}{AccentColor!20}

% --- Semantic aliases ---
\definecolor{HeadingColor}{named}{AccentColor}
\definecolor{LinkColor}{RGB}{38,52,139}
\definecolor{RuleColor}{named}{midgray}
\definecolor{TextColor}{named}{black}

% --- Page and text ---
\pagecolor{white}
\color{TextColor}

% --- Links ---
\hypersetup{
    colorlinks=true,
    linkcolor=LinkColor,
    citecolor=LinkColor,
    filecolor=LinkColor,
    urlcolor=LinkColor,
}

% ==============================================================
% 5. HEADER AND FOOTER
% ==============================================================
% Every field comes from cover.tex, which build_cover.py generates
% from the TOML files. Both journal logos sit in the running
% header at the same height, one at each margin; the footer
% carries the short citation on the left and the page count on
% the right.
\newcommand{\RBRHlogoheight}{8mm}
\newcommand{\HeaderLeft}{\includegraphics[height=\RBRHlogoheight]{\JournalLogoLeft}}
\newcommand{\HeaderRight}{\includegraphics[height=\RBRHlogoheight]{\JournalLogoRight}}
\newcommand{\FooterLeft}{\ArticleShortCitation}
\newcommand{\FooterCenter}{}
\newcommand{\FooterRight}{\thepage\;/\;\pageref{LastPage}}

\fancyhfoffset[L]{0mm}
\fancyhfoffset[R]{0mm}

\fancypagestyle{plain}{%
  \fancyhf{}
  \fancyhead[L]{\HeaderLeft}
  \fancyhead[R]{\HeaderRight}
  \renewcommand{\headrulewidth}{0.4pt}
  \fancyfoot[L]{\sffamily\scriptsize{\FooterLeft}}
  \fancyfoot[C]{\sffamily\scriptsize{\FooterCenter}}
  \fancyfoot[R]{\sffamily\scriptsize{\FooterRight}}
  \renewcommand{\footrulewidth}{0.2pt}
}

\pagestyle{fancy}
\fancyhf{}
\fancyhead[L]{\HeaderLeft}
\fancyhead[R]{\HeaderRight}
\renewcommand{\headrulewidth}{0.4pt}
\fancyfoot[L]{\sffamily\scriptsize{\FooterLeft}}
\fancyfoot[C]{\sffamily\scriptsize{\FooterCenter}}
\fancyfoot[R]{\sffamily\scriptsize{\FooterRight}}
\renewcommand{\footrulewidth}{0.2pt}

% ==============================================================
% 6. PAGE LAYOUT AND SPACING
% ==============================================================
\geometry{
 a4paper,
 left=20mm,
 right=20mm,
 top=20mm,
 includefoot,
 bottom=20mm,
 headheight=15mm,
 headsep=5mm,
 footskip=10mm,
}

\newcommand{\pagesectionvspace}{\vspace*{0mm}}

\setlength{\parindent}{0mm}
\setlength{\parskip}{2mm}

\setlength{\intextsep}{1mm}
\setlength{\textfloatsep}{2mm}
\setlength{\floatsep}{2mm}
\setlength{\abovecaptionskip}{1mm}
\setlength{\belowcaptionskip}{0mm}

\linespread{1}

\titlespacing*{\section}      {0mm}{4mm}{0mm}
\titlespacing*{\subsection}   {0mm}{2mm}{0mm}
\titlespacing*{\subsubsection}{0mm}{2mm}{0mm}
\titlespacing{\paragraph}{0mm}{2mm}{2mm}

% --------------------------------------------------------------
% Camera-ready sidebar column (left, narrow) + content column
% (right, wide). Width tuned for the DOI/highlights/correspondence
% sidebar used by main.tex; adjust \RBRHColumnRatio to change it.
\newcommand{\RBRHColumnRatio}{0.258}

% Lengths for the small-with-side-caption figure pattern (see
% contents.tex): a small, fixed-size image whose caption spans
% into the space the narrow sidebar column would otherwise
% occupy. Allocated once here; reused via \setlength (not
% \newlength) at the use site.
\newlength{\sidecapimgwidth}
\newlength{\sidecapgap}
\newlength{\sidecapwidth}
\newlength{\smallfigfixedwidth}

% ==============================================================
% 7. BIBLIOGRAPHY RESOURCES
% ==============================================================
\addbibresource{references.bib}
% Reference list one size down from the body text.
\renewcommand*{\bibfont}{\small}

% ==============================================================
% 8. CAPTIONS
% ==============================================================
\DeclareCaptionLabelSeparator{pipe}{\;|\;}

\captionsetup[figure]{
    font={footnotesize, sf, color=heavygray},
    labelfont={bf, color=AccentColor},
    labelsep=quad,
}

% No "width=" here. Inside paracol \textwidth is still the full
% page width, not the width of the column being typeset, so
% width=\textwidth made table captions overrun the text column and
% spill into the sidebar. Left unset, the caption takes the width
% of whatever contains it -- the text column for TableM, the full
% span for TableL -- which is right in both cases.
\captionsetup[table]{
    font={footnotesize, sf},
    labelfont={bf, color=AccentColor},
    labelsep=pipe,
}

\DeclareCaptionStyle{lstlisting}{
    font={footnotesize, sf},
    labelfont={bf, color=AccentColor},
    labelsep=pipe,
    justification=centering
}
\captionsetup[lstlisting]{style=lstlisting}

% All captions: justified, flush left, regardless of length.
% singlelinecheck=false disables caption's default of auto-centering
% any caption that happens to fit on a single line.
\captionsetup{justification=justified,singlelinecheck=false,labelfont=bf}

% ==============================================================
% 9. SECTION TITLES
% ==============================================================
\titleformat{\section}
  {\bfseries\sffamily\Large\color{HeadingColor}}
  {\makebox[8mm][l]{\thesection}}{2mm}{}

\titleformat{\subsection}
  {\bfseries\sffamily\large\color{HeadingColor}}
  {\makebox[8mm][l]{\thesubsection}}{2mm}{}

\titleformat{\subsubsection}
  {\sffamily\normalsize\color{HeadingColor}}
  {\makebox[10mm][l]{\thesubsubsection}}{4mm}{}

\titleformat{\paragraph}[runin]
  {\bfseries\sffamily\normalsize\color{HeadingColor}}
  {}{0mm}{}[.]

% Shared heading for auxiliary lists (References, List of Boxes,
% Glossary, ...) so they match the plain \section* styling used
% by \tableofcontents / \listoffigures / \listoftables, instead of
% falling back to each package's own internal heading.
\newcommand{\auxlisttitle}[1]{%
  \begingroup
    \titlespacing*{\section}{0mm}{6mm}{2mm}%
    \section*{#1}%
  \endgroup
  \phantomsection
  \addcontentsline{toc}{section}{#1}%
}

% ==============================================================
% 10. THEOREM-LIKE ENVIRONMENTS
% ==============================================================
\providecommand{\theoremname}{\trlang{Theorem}{Teorema}}
\providecommand{\lemmaname}{\trlang{Lemma}{Lema}}
\providecommand{\corollaryname}{\trlang{Corollary}{Corol\'ario}}
\providecommand{\definitionname}{\trlang{Definition}{Defini\c{c}\~ao}}
\providecommand{\examplename}{\trlang{Example}{Exemplo}}
\providecommand{\remarkname}{\trlang{Remark}{Observa\c{c}\~ao}}

\makeatletter
\renewcommand{\th@plain}{%
    \thm@headfont{\bfseries\sffamily\color{AccentColor}}%
    \itshape
}
\renewcommand{\th@definition}{%
    \thm@headfont{\bfseries\sffamily\color{AccentColor}}%
    \normalfont
}
\renewcommand{\th@remark}{%
    \thm@headfont{\itshape\sffamily\color{AccentColor}}%
    \normalfont
}
\makeatother

\newtheorem{theorem}{\theoremname}
\newtheorem{lemma}[theorem]{\lemmaname}
\newtheorem{corollary}[theorem]{\corollaryname}
\newtheorem{definition}{\definitionname}
\newtheorem{example}{\examplename}
\newtheorem*{remark}{\remarkname}

% ==============================================================
% 11. LISTINGS
% ==============================================================
\renewcommand{\lstlistingname}{\trListing}
\crefname{lstlisting}{\trListing}{\trListings}
\Crefname{lstlisting}{\trListing}{\trListings}

\lstdefinestyle{codeblock}{
    backgroundcolor=\color{black!5},
    commentstyle=\color{black!40}\itshape,
    keywordstyle=\color{purple}\bfseries,
    stringstyle=\color{teal},
    numberstyle=\tiny\color{black!40},
    basicstyle=\ttfamily\footnotesize,
    breaklines=true,
    breakatwhitespace=true,
    postbreak=\mbox{\textcolor{gray}{$\hookrightarrow$}\space},
    numbersep=8pt,
    tabsize=4,
    keepspaces=true,
    frame=none,
    captionpos=b,
    showspaces=false,
    showstringspaces=false,
    showtabs=false,
    numbers=left,
}
\lstset{style=codeblock, language=Python}
\lstset{basicstyle=\ttfamily\footnotesize}
\renewcommand{\ttdefault}{SourceCodePro-TLF}
\renewcommand{\familydefault}{\rmdefault}

% ==============================================================
% 12. BOXES
% ==============================================================
\tcbset{
    boxbasestyle/.style={
        enhanced,
        breakable,
        colback      = AccentColorVeryLight,
        colbacktitle = AccentColorLight,
        coltitle     = AccentColor,
        colframe     = AccentColor,
        boxrule      = 0.5pt,
        arc          = 2pt,
        fonttitle    = \bfseries\sffamily\small,
        fontupper    = \sffamily\footnotesize,
        boxsep       = 0mm,
        top          = 2mm,
        bottom       = 2mm,
        left         = 2mm,
        right        = 2mm,
        toptitle     = 1mm,
        bottomtitle  = 1mm,
    }
}

% boxstory — numbered, breakable, document-wide counter.
% Usage: \begin{boxstory}{Title} ... \end{boxstory}
\newtcolorbox[
    auto counter,
    list inside = lob,
]{boxstory}[2][]{
    boxbasestyle,
    title = {\sffamily Box~\thetcbcounter\quad #2},
    #1,
}

\newcommand{\notitleheading}[1]{}
\newcommand{\listofboxes}{%
    \auxlisttitle{\trListOfBoxes}%
    \tcblistof[\notitleheading]{lob}{\trListOfBoxes}%
}

% ==============================================================
% 13. ADMONITIONS
% ==============================================================
% Color-coded unnumbered admonition boxes. Icons require fontawesome5.
\newtcolorbox{amdnote}{
    boxbasestyle,
    colback = cyan!5, colbacktitle = cyan!50!black, colframe = cyan!50!black,
    coltitle = white, title = {\faInfoCircle\quad \trAdmNote},
}
\newtcolorbox{amdtip}{
    boxbasestyle,
    colback = green!5!white, colbacktitle = green!40!black, colframe = green!40!black,
    coltitle = white, title = {\faLightbulb\quad \trAdmTip},
}
\newtcolorbox{amdimportant}{
    boxbasestyle,
    colback = violet!5, colbacktitle = violet!60!black, colframe = violet!60!black,
    coltitle = white, title = {\faExclamationCircle\quad \trAdmImportant},
}
\newtcolorbox{amdwarning}{
    boxbasestyle,
    colback = orange!5, colbacktitle = orange!70!black, colframe = orange!70!black,
    coltitle = white, title = {\faExclamationTriangle\quad \trAdmWarning},
}
\newtcolorbox{amdcaution}{
    boxbasestyle,
    colback = red!5, colbacktitle = red!70!black, colframe = red!70!black,
    coltitle = white, title = {\faBan\quad \trAdmCaution},
}

% ==============================================================
% 14. DRAFT WATERMARK
% ==============================================================
% main.tex sets \IsDraft to "true" to stamp every page "DRAFT" /
% "MINUTA" while the manuscript is a work in progress.
\providecommand{\IsDraft}{false}
\ifdefstring{\IsDraft}{true}{
    \usepackage{draftwatermark}
    \SetWatermarkText{\trDraft}
    \SetWatermarkScale{1}
    \SetWatermarkColor[gray]{0.8}
    \SetWatermarkAngle{30}
}{}

% ==============================================================
% 15. COVER
% ==============================================================
% Everything here reads macros defined in cover.tex, which
% build_cover.py generates from journal.toml, article.toml and
% authors.toml. Nothing on the cover is written by hand, which is
% why contents.tex can begin straight at the first section.
%
% \PrintMasthead fills the wide column: journal name, title,
% authors, affiliations, then both abstracts. \PrintSidebar fills
% the narrow one, in the order RBRH prints it.

% The author and affiliation blocks are long runs of proper nouns
% and institution names, which offer TeX very few legal break
% points. On the default settings the author list overflowed the
% column by about 10pt.
%
% Justified was tried first and does fix the overflow, with
% \tolerance raised and \emergencystretch added -- but it pays for
% it by hyphenating personal names across lines ("An-derson",
% "Fer-nando") and by leaving one very loose line. Breaking a
% surname is not acceptable in a masthead, so these blocks are set
% ragged right instead, with hyphenation switched off outright.
% The right edge is uneven; nothing overflows and no name breaks.
%
% To go back to justified, swap \raggedright and the two penalties
% for: \tolerance=4000 \emergencystretch=2em
%
% These are paragraph-breaking parameters, so they only take
% effect if the paragraph ends while they are still in scope --
% hence the \par inside each group below.
\newcommand{\RBRHnameblock}{%
  \raggedright
  \hyphenpenalty=10000
  \exhyphenpenalty=10000
}

% --- Wide column ---
\newcommand{\PrintMasthead}{%
  {\noindent \sffamily \Large\textbf{\JournalName}}\\
  {\noindent \large \textit{\JournalNameEn}}\\[3mm]
  \begin{center}
    \raggedright {\Large \sffamily \textbf{\color{AccentColor}\ArticleTitle}}\\[3mm]
  \end{center}
  {\RBRHnameblock\noindent\AuthorBlock\par}
  \vspace{2mm}
  % Affiliations in the same grey as the captions, so they read as
  % secondary to the author names without fading out like the
  % licence notice does.
  {\footnotesize\color{heavygray}\RBRHnameblock\noindent\AffiliationBlock\par}
  \vspace{2mm}
  \hrule
  \vspace{2mm}
  \PrintAbstracts
}

% Both languages are printed whatever \DocLanguage says: RBRH
% carries an English abstract and a Portuguese resumo regardless
% of the language the body is written in.
% The \small applies to the heading words only -- "Abstract",
% "Keywords:" and their Portuguese counterparts, which were
% sitting at full body size. The abstract text itself stays at
% \footnotesize.
\newcommand{\PrintAbstracts}{%
  \begin{spacing}{0.8}
  \small
  {\noindent \sffamily\textbf{\color{AccentColor}Abstract}\quad{\footnotesize\color{TextColor} \AbstractEn}}\\[1mm]
  {\noindent \sffamily\textbf{\color{AccentColor}Keywords:}\quad{\footnotesize\color{TextColor} \KeywordsEn}}\\[3mm]
  {\noindent \sffamily\textbf{\color{AccentColor}Resumo}\quad{\footnotesize\color{TextColor} \AbstractPt}}\\[1mm]
  {\noindent \sffamily\textbf{\color{AccentColor}Palavras-chave:}\quad{\footnotesize\color{TextColor} \KeywordsPt}}
  \end{spacing}
}

% --- Narrow column ---
% One helper per block, so the running order is readable and can
% be reshuffled by editing \PrintSidebar alone.
% The accent marks the structure of the column -- the heading of
% each block -- while the content under it stays in body colour.
\newcommand{\RBRHcoverblock}[2]{%
  {\color{AccentColor}\textbf{#1}}\\[1mm]
  #2
  \\[4mm]
}

\newcommand{\PrintSidebar}{%
  {
  \sffamily
  \scriptsize
  \raggedright
  % licence block at the foot overrides it with grey.
  \textbf{\MakeUppercase{\ArticleType}}\\
  \ifdefempty{\ArticleDOI}{}{DOI\;\ArticleDOI\\[4mm]}
  \RBRHcoverblock{\trHighlights:}{\ArticleHighlights}
  \RBRHcoverblock{\trCitation:}{\ArticleCitation}
  \RBRHcoverblock{\trCorrespondence:}{\CorrespondingBlock}
  \RBRHcoverblock{\trCompetingInterests:}{\CompetingInterests}
  \RBRHcoverblock{\trAcknowledgements:}{\Acknowledgements}
  \RBRHcoverblock{\trEditors:}{%
    \trEditorInChief:\\ \EditorInChief\\[1mm]
    \trAssociatedEditor:\\ \AssociatedEditor}
  \RBRHcoverblock{\trDataAvailability:}{\DataAvailability}
  \RBRHcoverblock{\trPublicationDates:}{%
    \trPublished\ \DatePublished\\
    \trIssue\ \ArticleIssue}
  \RBRHcoverblock{\trHistory:}{%
    \trReceived\ \DateReceived\\
    \trReviewed\ \DateReviewed\\
    \trAccepted\ \DateAccepted}
  % Push the licence to the foot of the column: everything above
  % sets at its natural height and this absorbs the slack, so the
  % notice sits on the baseline of the page however much or little
  % metadata precedes it.
  \vfill
  \PrintLicense
  }
}

% CC-BY notice, set in grey so it reads as a footer to the
% sidebar rather than as content. The badge is a vector PDF
% converted from assets/cc-by.svg.
\newcommand{\PrintLicense}{%
  {\color{midgray}
   \includegraphics[height=3.2mm]{\LicenseBadge}\\[1mm]
   \LicenseCopyright\\
   \LicenseStatement\par}
}

% ==============================================================
% 16. APPENDIX RENUMBERING
% ==============================================================
% Call \BeginAppendix in contents.tex to switch into lettered,
% restarted numbering (A, A.1, Figure A.1, ...) for supplementary
% material.
\newcommand{\BeginAppendix}{%
  \appendix
  \renewcommand{\thesection}{\Alph{section}}
  \renewcommand{\thesubsection}{\Alph{section}.\arabic{subsection}}
  \renewcommand{\thesubsubsection}{\Alph{section}.\arabic{subsection}.\arabic{subsubsection}}
  \setcounter{section}{0}
  \setcounter{equation}{0}
  \setcounter{figure}{0}
  \setcounter{table}{0}
  \setcounter{lstlisting}{0}
  \renewcommand{\theequation}{\Alph{section}.\arabic{equation}}
  \renewcommand{\thefigure}{\Alph{section}.\arabic{figure}}
  \renewcommand{\thetable}{\Alph{section}.\arabic{table}}
  \renewcommand{\thelstlisting}{\Alph{section}.\arabic{lstlisting}}
  \pagesectionvspace
  \auxlisttitle{\trAppendices}
  \label{sec:appendix}
  \vspace{10mm}
  \small
  % Headings one level down too. The body is already \small, which
  % left \Large section titles looking oversized against it.
  \titleformat{\section}
    {\bfseries\sffamily\large\color{HeadingColor}}
    {\makebox[8mm][l]{\thesection}}{2mm}{}
  \titleformat{\subsection}
    {\bfseries\sffamily\normalsize\color{HeadingColor}}
    {\makebox[8mm][l]{\thesubsection}}{2mm}{}
  \titleformat{\subsubsection}
    {\sffamily\small\color{HeadingColor}}
    {\makebox[10mm][l]{\thesubsubsection}}{4mm}{}
}

% ==============================================================
% 17. SHARED COUNTERS ACROSS PARACOL COLUMNS
% ==============================================================
% paracol makes every counter column-LOCAL by default -- its
% internal list of global counters ships as just {page}. Each
% column therefore keeps its own `figure`, `table` and `equation`.
% A figure written inside a column and one written in a
% \switchcolumn[1]*[...] spanning block are stepped in different
% counter states, so the spanning one restarts at 1 and the
% numbering appears to repeat. `figure` is shared instead: the two
% columns here are a metadata sidebar plus one article body, not
% parallel independent texts.
%
% Name the counter explicitly -- do NOT use \globalcounter*, and
% do NOT add `table` here. Both were tried and both break this
% document with "Missing \endcsname inserted" at the first
% \section (the starred form because it sweeps in every counter,
% including `table`). Bisected against this exact package set:
%   figure   -> safe, and is what fixes the numbering bug
%   equation -> safe, but unused; add only if an equation ever
%               goes inside a \switchcolumn[1]*[...] span
%   table    -> BREAKS the build, do not enable
% `footnote` is left alone too: paracol manages its scope itself.
\globalcounter{figure}

% ==============================================================
% 18. AUTHOR-FACING FIGURE AND TABLE MACROS
% ==============================================================
% So contents.tex never carries layout plumbing:
%
%   \FigureM[pos]{name}{caption}{label}        text width
%   \FigureS[pos][frac]{name}{caption}{label}  small -- frac of the
%                                              text width, default
%                                              0.67 -- and the only
%                                              size whose caption
%                                              sits beside the
%                                              image, spanning into
%                                              the sidebar column
%   \FigureL[pos]{name}{caption}{label}        full width, margin to
%                                              margin, matching the
%                                              width of the header
%                                              rule
%
%   \begin{TableM}[pos]{caption}{label} ... \end{TableM}  text width
%   \begin{TableL}[pos]{caption}{label} ... \end{TableL}  full width
%
% [pos] is the usual LaTeX float placement specifier and defaults
% to t (top). The table environments take the tabular material as
% their body and put the caption above it, per table convention;
% figures caption below.
%
% <name> is the bare file name: no folder, no extension. The file
% is resolved inside figures/, preferring .pdf, then .png, then
% .jpeg, then .jpg. A name that matches nothing degrades to a
% placeholder box rather than failing the build.
%
% M/S/L map onto real floats -- figure/table for M, and the
% double-column figure*/table* for S and L, which is what makes
% them span the sidebar. paracol supports both (it carries its own
% \pcol@floatplacement and spanning-float machinery), and the
% figure counter is shared across columns by section 19 above, so
% numbering follows source order. \label goes straight after
% \caption inside the float, where \@currentlabel is correct.

\newif\ifRBRHfigfound
\newcommand{\RBRHfigpath}{}

% Resolve figures/<name> against the supported extensions, in
% order of preference.
\newcommand{\RBRHfindfigure}[1]{%
  \IfFileExists{figures/#1.pdf}%
    {\RBRHfigfoundtrue \def\RBRHfigpath{figures/#1.pdf}}%
    {\IfFileExists{figures/#1.png}%
      {\RBRHfigfoundtrue \def\RBRHfigpath{figures/#1.png}}%
      {\IfFileExists{figures/#1.jpeg}%
        {\RBRHfigfoundtrue \def\RBRHfigpath{figures/#1.jpeg}}%
        {\IfFileExists{figures/#1.jpg}%
          {\RBRHfigfoundtrue \def\RBRHfigpath{figures/#1.jpg}}%
          {\RBRHfigfoundfalse \def\RBRHfigpath{}}}}}%
}

% Placeholder for a missing file: gray box in golden ratio
% (height = width / 1.618) carrying a red notice. \detokenize
% keeps names containing "_" printable.
\newlength{\RBRHfigwidth}
\newlength{\RBRHfigheight}
\newcommand{\RBRHmissingfigure}[2]{%
  \setlength{\RBRHfigwidth}{#1}%
  \setlength{\RBRHfigheight}{\dimexpr\RBRHfigwidth*1000/1618\relax}%
  \begingroup
    \setlength{\fboxsep}{0pt}%
    \colorbox{lightgray}{%
      \begin{minipage}[c][\RBRHfigheight][c]{\RBRHfigwidth}%
        \centering
        \textcolor{red}{\sffamily\small
          Figure \texttt{\detokenize{#2}} not found}%
      \end{minipage}%
    }%
  \endgroup
}

% The image at #1 wide, or the placeholder if it is missing.
\newcommand{\RBRHfigbody}[2]{%
  \ifRBRHfigfound
    \includegraphics[width=#1]{\RBRHfigpath}%
  \else
    \RBRHmissingfigure{#1}{#2}%
  \fi
}

% --- M: text width ---
\NewDocumentCommand{\FigureM}{O{t} m m m}{%
  \begin{figure}[#1]
    \centering
    \RBRHfindfigure{#2}%
    \RBRHfigbody{\linewidth}{#2}%
    \caption{#3}\label{#4}%
  \end{figure}%
}

% --- L: full width, margin to margin ---
\NewDocumentCommand{\FigureL}{O{t} m m m}{%
  \begin{figure*}[#1]
    \centering
    \RBRHfindfigure{#2}%
    \RBRHfigbody{\textwidth}{#2}%
    \caption{#3}\label{#4}%
  \end{figure*}%
}

% --- S: small, caption beside the image, spanning the sidebar ---
% The image width is captured BEFORE the float opens, while
% \linewidth still means the text column: inside a double-column
% float LaTeX widens it to \textwidth. The float body is typeset
% here at the call site, so the captured value is the right one.
\NewDocumentCommand{\FigureS}{O{t} O{0.67} m m m}{%
  \setlength{\smallfigfixedwidth}{#2\linewidth}%
  \begin{figure*}[#1]
    \RBRHfindfigure{#3}%
    \setlength{\sidecapimgwidth}{\smallfigfixedwidth}%
    \setlength{\sidecapgap}{6mm}%
    \setlength{\sidecapwidth}{\textwidth}%
    \addtolength{\sidecapwidth}{-\sidecapimgwidth}%
    \addtolength{\sidecapwidth}{-\sidecapgap}%
    \noindent
    % \vspace{0pt} as the first thing in each minipage is what makes
    % [t] align the TOPS. Without it, [t] aligns first baselines, and
    % a minipage holding just a graphic has its baseline at the
    % graphic's bottom -- which dropped the caption to the foot of
    % the image instead of level with its top. \abovecaptionskip is
    % zeroed for the same reason: it is meant to separate a figure
    % from a caption printed below it, not beside it.
    \begin{minipage}[t]{\sidecapwidth}
      \vspace{0pt}%
      \setlength{\abovecaptionskip}{0pt}%
      \caption{#4}\label{#5}%
    \end{minipage}%
    \hspace{\sidecapgap}%
    \begin{minipage}[t]{\sidecapimgwidth}
      \vspace{0pt}%
      \RBRHfigbody{\linewidth}{#3}%
    \end{minipage}%
  \end{figure*}%
}

% --- Tables: caption above the body, per table convention ---
% The body is wrapped in threeparttable, so the caption and any
% notes take the width of the tabular rather than running wider,
% and authors can add a notes block without extra scaffolding:
%
%   \begin{TableM}{caption}{label}
%     \begin{tabular}{ll} ... \end{tabular}
%     \begin{tablenotes}
%       \item Note text.
%     \end{tablenotes}
%   \end{TableM}
\NewDocumentEnvironment{TableM}{O{t} m m}{%
  \begin{table}[#1]
    \centering
    \begin{threeparttable}
      \caption{#2}\label{#3}%
}{%
    \end{threeparttable}
  \end{table}%
}

\NewDocumentEnvironment{TableL}{O{t} m m}{%
  \begin{table*}[#1]
    \centering
    \begin{threeparttable}
      \caption{#2}\label{#3}%
}{%
    \end{threeparttable}
  \end{table*}%
}
;
% this \providecommand only guards against main.tex forgetting to,
% and matches the template's own default: the document is English.
\providecommand{\DocLanguage}{english}

% ==============================================================
% 1. PACKAGES
% ==============================================================

% --------------------------------------------------------------
% Fonts
% Times (serif, body text + math) + Helvetica (sans, headings/
% captions) + Source Code Pro (mono, verbatim/lstlisting).
% All bundled with standard TeX Live; compatible with pdfLaTeX.
\usepackage[T1]{fontenc}
\usepackage[utf8]{inputenc}
\usepackage{mathptmx}                    % Times, incl. matching math glyphs
\usepackage[scaled=0.92]{helvet}         % Helvetica
% Source Code Pro has a tall x-height and reads a size larger than
% Times beside it, so it is scaled down to match rather than left
% to dominate every \texttt and listing.
\usepackage[scaled=0.85]{sourcecodepro}

% --------------------------------------------------------------
% Math and symbols
\usepackage{amsmath}
\usepackage{amssymb}
\usepackage{amsthm}
\usepackage{textcomp}

% --------------------------------------------------------------
% Language, references and datetime
% cleveref and hyperref load last, after babel, per their
% required order: babel -> cleveref -> hyperref.
% Both languages are loaded, and \DocLanguage picks the MAIN one.
% "main=" is what actually decides the document language: babel
% otherwise takes the last language in the option list as main and
% re-asserts it at \begin{document}, overriding any preamble
% \selectlanguage. That is what kept captions reading "Figura" and
% the footer date in Portuguese under \DocLanguage=english.
% The option list has to be built by \edef first, because package
% options are not expanded when \usepackage reads them.
\edef\RBRHbabelopts{main=\DocLanguage,english,portuguese}
\expandafter\PassOptionsToPackage\expandafter{\RBRHbabelopts}{babel}
\usepackage{babel}
\usepackage[useregional]{datetime2}
\usepackage{csquotes}
\usepackage{hyperref}
\usepackage{cleveref}

% --------------------------------------------------------------
% Colors
\usepackage[dvipsnames, table]{xcolor}

% --------------------------------------------------------------
% Layout
\usepackage{geometry}
\usepackage{setspace}
\usepackage{titlesec}
\usepackage{fancyhdr}
\usepackage{lastpage}
\usepackage{paracol}                     % the two-column "sidebar + content" skeleton

% --------------------------------------------------------------
% Bibliography
% APA style with author-year citations, biber backend. Authors
% cite with \textcite{} in prose and \parencite{} in parentheses.
\usepackage[backend=biber, style=apa, citestyle=authoryear]{biblatex}

% --------------------------------------------------------------
% Sections
\usepackage{etoolbox}

% --------------------------------------------------------------
% Text
\usepackage{lipsum}                      % placeholder text only; not needed once authors write

% --------------------------------------------------------------
% Captions
\usepackage{caption}

% --------------------------------------------------------------
% Figures
\usepackage{graphicx}

% --------------------------------------------------------------
% Tables
\usepackage{array}
\usepackage{booktabs}
\usepackage{threeparttable}
% tabularx so a table can be told to occupy exactly the width it
% is given. threeparttable sizes the caption and the notes to the
% width of the tabular, so a table that fills \linewidth is what
% keeps both flush with the text column instead of ending short.
\usepackage{tabularx}
% Left-aligned flexible column: the X column, minus the centring.
\newcolumntype{L}{>{\raggedright\arraybackslash}X}

% --------------------------------------------------------------
% Boxes
\usepackage{tcolorbox}
\tcbuselibrary{skins, breakable}

% --------------------------------------------------------------
% Listings
\usepackage{listings}

% --------------------------------------------------------------
% Misc
\usepackage{fontawesome5}                % admonition icons
% ORCID mark for the author list. Draws the official logo as
% vector paths rather than embedding an SVG, which pdfLaTeX
% cannot do without --shell-escape. It scales the mark to the
% current font's X-height, so it sits inside the line and cannot
% change the leading, and it makes the icon itself the hyperlink.
\usepackage{orcidlink}

% ==============================================================
% 2. LANGUAGE SWITCH
% ==============================================================
% Nothing to do here: the document language is set by babel's
% "main=" option in section 1, which is the only thing that
% survives \begin{document}. It drives hyphenation, babel-aware
% dates (datetime2 [useregional]), quoting (csquotes) and the
% caption words ("Figure"/"Figura", "Table"/"Tabela").
%
% A preamble \selectlanguage used to sit here and did nothing:
% babel overrode it on entering the document body.

% ==============================================================
% 3. TRANSLATIONS
% ==============================================================
% Journal-facing labels, in Portuguese or English depending on
% \DocLanguage. Authors reference these by name (e.g. \trAbstract)
% rather than hardcoding either language, so the whole document
% switches with one macro.
% \ifx rather than etoolbox's \ifdefstring: \ifx is expandable, so
% \trlang survives \MakeUppercase. With \ifdefstring the conditional
% gets mangled during uppercasing and the wrong branch wins -- that
% is what made the sidebar print "ESTUDO DE CASO" under english.
%
% Both sides are normalised into plain macros first. \ifx compares
% meanings, and \newcommand makes \DocLanguage \long while a bare
% \def would not -- comparing them directly fails on that prefix
% alone and silently selects Portuguese. \edef flattens whatever
% the author wrote into the same shape as \RBRHenglish.
%
% No "@" in these names: this is a .tex file, not a .sty, so @ is
% not a letter here and \RBRHxx@yy would parse as \RBRHxx followed
% by stray characters.
\edef\RBRHlang{\DocLanguage}
\def\RBRHenglish{english}
\newcommand{\trlang}[2]{\ifx\RBRHlang\RBRHenglish #1\else #2\fi}

% --- Front matter ---
\newcommand{\trAbstract}{\trlang{Abstract}{Resumo}}
\newcommand{\trKeywords}{\trlang{Keywords}{Palavras-chave}}
\newcommand{\trListOfBoxes}{\trlang{Boxes}{Destaques}}
\newcommand{\trHighlights}{\trlang{Highlights}{Destaques do estudo}}
\newcommand{\trCorrespondence}{\trlang{Correspondence to}{Correspond\^encia para}}
\newcommand{\trCitation}{\trlang{Citation}{Como citar}}

% --- Cover sidebar ---
\newcommand{\trEditors}{\trlang{Editors}{Editores}}
\newcommand{\trEditorInChief}{\trlang{Editor-in-Chief}{Editor-chefe}}
\newcommand{\trAssociatedEditor}{\trlang{Associated Editor}{Editor associado}}
\newcommand{\trCompetingInterests}{\trlang{Competing interests}{Conflito de interesses}}
\newcommand{\trDataAvailability}{\trlang{Data availability}{Disponibilidade dos dados}}
\newcommand{\trPublicationDates}{\trlang{Publication dates}{Datas de publica\c{c}\~ao}}
\newcommand{\trPublished}{\trlang{Published}{Publicado em}}
\newcommand{\trIssue}{\trlang{Issue}{Edi\c{c}\~ao}}
\newcommand{\trHistory}{\trlang{History}{Hist\'orico}}
\newcommand{\trReceived}{\trlang{Received}{Recebido em}}
\newcommand{\trReviewed}{\trlang{Reviewed}{Revisado em}}
\newcommand{\trAccepted}{\trlang{Accepted}{Aceito em}}

% --- Back matter ---
\newcommand{\trReferences}{\trlang{References}{Refer\^encias}}
\newcommand{\trAppendix}{\trlang{Appendix}{Ap\^endice}}
\newcommand{\trAppendices}{\trlang{Appendices}{Ap\^endices}}
\newcommand{\trAcknowledgements}{\trlang{Acknowledgements}{Agradecimentos}}

% --- Items ---
\newcommand{\trListing}{\trlang{Code}{C\'odigo}}
\newcommand{\trListings}{\trlang{Code}{C\'odigos}}
\newcommand{\trDraft}{\trlang{DRAFT}{MINUTA}}

% --- Admonitions ---
\newcommand{\trAdmNote}{\trlang{Note}{Nota}}
\newcommand{\trAdmTip}{\trlang{Tip}{Dica}}
\newcommand{\trAdmImportant}{\trlang{Important}{Importante}}
\newcommand{\trAdmWarning}{\trlang{Warning}{Alerta}}
\newcommand{\trAdmCaution}{\trlang{Caution}{Cuidado}}

% ==============================================================
% 4. COLORS
% ==============================================================
% Override this section to apply a different brand/theme.

% --- Grayscale utilities ---
\definecolor{verylightgray}{gray}{0.97}
\definecolor{lightgray}{gray}{0.85}
\definecolor{midgray}{gray}{0.50}
\definecolor{heavygray}{gray}{0.3}

% --- Brand / accent ---
% The journal's blue is the accent colour: it carries the article
% title, the section headings, the abstract headings and the
% metadata column. The masthead itself is set in black, so the
% blue marks structure rather than branding.
%
% Everything already keyed to AccentColor follows from here --
% caption labels, theorem heads and the tcolorbox palette, whose
% tints derive from it.
\definecolor{AccentColor}{RGB}{38,52,139}
\colorlet{AccentColorVeryLight}{AccentColor!5}
\colorlet{AccentColorLight}{AccentColor!20}

% --- Semantic aliases ---
\definecolor{HeadingColor}{named}{AccentColor}
\definecolor{LinkColor}{RGB}{38,52,139}
\definecolor{RuleColor}{named}{midgray}
\definecolor{TextColor}{named}{black}

% --- Page and text ---
\pagecolor{white}
\color{TextColor}

% --- Links ---
\hypersetup{
    colorlinks=true,
    linkcolor=LinkColor,
    citecolor=LinkColor,
    filecolor=LinkColor,
    urlcolor=LinkColor,
}

% ==============================================================
% 5. HEADER AND FOOTER
% ==============================================================
% Every field comes from cover.tex, which build_cover.py generates
% from the TOML files. Both journal logos sit in the running
% header at the same height, one at each margin; the footer
% carries the short citation on the left and the page count on
% the right.
\newcommand{\RBRHlogoheight}{8mm}
\newcommand{\HeaderLeft}{\includegraphics[height=\RBRHlogoheight]{\JournalLogoLeft}}
\newcommand{\HeaderRight}{\includegraphics[height=\RBRHlogoheight]{\JournalLogoRight}}
\newcommand{\FooterLeft}{\ArticleShortCitation}
\newcommand{\FooterCenter}{}
\newcommand{\FooterRight}{\thepage\;/\;\pageref{LastPage}}

\fancyhfoffset[L]{0mm}
\fancyhfoffset[R]{0mm}

\fancypagestyle{plain}{%
  \fancyhf{}
  \fancyhead[L]{\HeaderLeft}
  \fancyhead[R]{\HeaderRight}
  \renewcommand{\headrulewidth}{0.4pt}
  \fancyfoot[L]{\sffamily\scriptsize{\FooterLeft}}
  \fancyfoot[C]{\sffamily\scriptsize{\FooterCenter}}
  \fancyfoot[R]{\sffamily\scriptsize{\FooterRight}}
  \renewcommand{\footrulewidth}{0.2pt}
}

\pagestyle{fancy}
\fancyhf{}
\fancyhead[L]{\HeaderLeft}
\fancyhead[R]{\HeaderRight}
\renewcommand{\headrulewidth}{0.4pt}
\fancyfoot[L]{\sffamily\scriptsize{\FooterLeft}}
\fancyfoot[C]{\sffamily\scriptsize{\FooterCenter}}
\fancyfoot[R]{\sffamily\scriptsize{\FooterRight}}
\renewcommand{\footrulewidth}{0.2pt}

% ==============================================================
% 6. PAGE LAYOUT AND SPACING
% ==============================================================
\geometry{
 a4paper,
 left=20mm,
 right=20mm,
 top=20mm,
 includefoot,
 bottom=20mm,
 headheight=15mm,
 headsep=5mm,
 footskip=10mm,
}

\newcommand{\pagesectionvspace}{\vspace*{0mm}}

\setlength{\parindent}{0mm}
\setlength{\parskip}{2mm}

\setlength{\intextsep}{1mm}
\setlength{\textfloatsep}{2mm}
\setlength{\floatsep}{2mm}
\setlength{\abovecaptionskip}{1mm}
\setlength{\belowcaptionskip}{0mm}

\linespread{1}

\titlespacing*{\section}      {0mm}{4mm}{0mm}
\titlespacing*{\subsection}   {0mm}{2mm}{0mm}
\titlespacing*{\subsubsection}{0mm}{2mm}{0mm}
\titlespacing{\paragraph}{0mm}{2mm}{2mm}

% --------------------------------------------------------------
% Camera-ready sidebar column (left, narrow) + content column
% (right, wide). Width tuned for the DOI/highlights/correspondence
% sidebar used by main.tex; adjust \RBRHColumnRatio to change it.
\newcommand{\RBRHColumnRatio}{0.258}

% Lengths for the small-with-side-caption figure pattern (see
% contents.tex): a small, fixed-size image whose caption spans
% into the space the narrow sidebar column would otherwise
% occupy. Allocated once here; reused via \setlength (not
% \newlength) at the use site.
\newlength{\sidecapimgwidth}
\newlength{\sidecapgap}
\newlength{\sidecapwidth}
\newlength{\smallfigfixedwidth}

% ==============================================================
% 7. BIBLIOGRAPHY RESOURCES
% ==============================================================
\addbibresource{references.bib}
% Reference list one size down from the body text.
\renewcommand*{\bibfont}{\small}

% ==============================================================
% 8. CAPTIONS
% ==============================================================
\DeclareCaptionLabelSeparator{pipe}{\;|\;}

\captionsetup[figure]{
    font={footnotesize, sf, color=heavygray},
    labelfont={bf, color=AccentColor},
    labelsep=quad,
}

% No "width=" here. Inside paracol \textwidth is still the full
% page width, not the width of the column being typeset, so
% width=\textwidth made table captions overrun the text column and
% spill into the sidebar. Left unset, the caption takes the width
% of whatever contains it -- the text column for TableM, the full
% span for TableL -- which is right in both cases.
\captionsetup[table]{
    font={footnotesize, sf},
    labelfont={bf, color=AccentColor},
    labelsep=pipe,
}

\DeclareCaptionStyle{lstlisting}{
    font={footnotesize, sf},
    labelfont={bf, color=AccentColor},
    labelsep=pipe,
    justification=centering
}
\captionsetup[lstlisting]{style=lstlisting}

% All captions: justified, flush left, regardless of length.
% singlelinecheck=false disables caption's default of auto-centering
% any caption that happens to fit on a single line.
\captionsetup{justification=justified,singlelinecheck=false,labelfont=bf}

% ==============================================================
% 9. SECTION TITLES
% ==============================================================
\titleformat{\section}
  {\bfseries\sffamily\Large\color{HeadingColor}}
  {\makebox[8mm][l]{\thesection}}{2mm}{}

\titleformat{\subsection}
  {\bfseries\sffamily\large\color{HeadingColor}}
  {\makebox[8mm][l]{\thesubsection}}{2mm}{}

\titleformat{\subsubsection}
  {\sffamily\normalsize\color{HeadingColor}}
  {\makebox[10mm][l]{\thesubsubsection}}{4mm}{}

\titleformat{\paragraph}[runin]
  {\bfseries\sffamily\normalsize\color{HeadingColor}}
  {}{0mm}{}[.]

% Shared heading for auxiliary lists (References, List of Boxes,
% Glossary, ...) so they match the plain \section* styling used
% by \tableofcontents / \listoffigures / \listoftables, instead of
% falling back to each package's own internal heading.
\newcommand{\auxlisttitle}[1]{%
  \begingroup
    \titlespacing*{\section}{0mm}{6mm}{2mm}%
    \section*{#1}%
  \endgroup
  \phantomsection
  \addcontentsline{toc}{section}{#1}%
}

% ==============================================================
% 10. THEOREM-LIKE ENVIRONMENTS
% ==============================================================
\providecommand{\theoremname}{\trlang{Theorem}{Teorema}}
\providecommand{\lemmaname}{\trlang{Lemma}{Lema}}
\providecommand{\corollaryname}{\trlang{Corollary}{Corol\'ario}}
\providecommand{\definitionname}{\trlang{Definition}{Defini\c{c}\~ao}}
\providecommand{\examplename}{\trlang{Example}{Exemplo}}
\providecommand{\remarkname}{\trlang{Remark}{Observa\c{c}\~ao}}

\makeatletter
\renewcommand{\th@plain}{%
    \thm@headfont{\bfseries\sffamily\color{AccentColor}}%
    \itshape
}
\renewcommand{\th@definition}{%
    \thm@headfont{\bfseries\sffamily\color{AccentColor}}%
    \normalfont
}
\renewcommand{\th@remark}{%
    \thm@headfont{\itshape\sffamily\color{AccentColor}}%
    \normalfont
}
\makeatother

\newtheorem{theorem}{\theoremname}
\newtheorem{lemma}[theorem]{\lemmaname}
\newtheorem{corollary}[theorem]{\corollaryname}
\newtheorem{definition}{\definitionname}
\newtheorem{example}{\examplename}
\newtheorem*{remark}{\remarkname}

% ==============================================================
% 11. LISTINGS
% ==============================================================
\renewcommand{\lstlistingname}{\trListing}
\crefname{lstlisting}{\trListing}{\trListings}
\Crefname{lstlisting}{\trListing}{\trListings}

\lstdefinestyle{codeblock}{
    backgroundcolor=\color{black!5},
    commentstyle=\color{black!40}\itshape,
    keywordstyle=\color{purple}\bfseries,
    stringstyle=\color{teal},
    numberstyle=\tiny\color{black!40},
    basicstyle=\ttfamily\footnotesize,
    breaklines=true,
    breakatwhitespace=true,
    postbreak=\mbox{\textcolor{gray}{$\hookrightarrow$}\space},
    numbersep=8pt,
    tabsize=4,
    keepspaces=true,
    frame=none,
    captionpos=b,
    showspaces=false,
    showstringspaces=false,
    showtabs=false,
    numbers=left,
}
\lstset{style=codeblock, language=Python}
\lstset{basicstyle=\ttfamily\footnotesize}
\renewcommand{\ttdefault}{SourceCodePro-TLF}
\renewcommand{\familydefault}{\rmdefault}

% ==============================================================
% 12. BOXES
% ==============================================================
\tcbset{
    boxbasestyle/.style={
        enhanced,
        breakable,
        colback      = AccentColorVeryLight,
        colbacktitle = AccentColorLight,
        coltitle     = AccentColor,
        colframe     = AccentColor,
        boxrule      = 0.5pt,
        arc          = 2pt,
        fonttitle    = \bfseries\sffamily\small,
        fontupper    = \sffamily\footnotesize,
        boxsep       = 0mm,
        top          = 2mm,
        bottom       = 2mm,
        left         = 2mm,
        right        = 2mm,
        toptitle     = 1mm,
        bottomtitle  = 1mm,
    }
}

% boxstory — numbered, breakable, document-wide counter.
% Usage: \begin{boxstory}{Title} ... \end{boxstory}
\newtcolorbox[
    auto counter,
    list inside = lob,
]{boxstory}[2][]{
    boxbasestyle,
    title = {\sffamily Box~\thetcbcounter\quad #2},
    #1,
}

\newcommand{\notitleheading}[1]{}
\newcommand{\listofboxes}{%
    \auxlisttitle{\trListOfBoxes}%
    \tcblistof[\notitleheading]{lob}{\trListOfBoxes}%
}

% ==============================================================
% 13. ADMONITIONS
% ==============================================================
% Color-coded unnumbered admonition boxes. Icons require fontawesome5.
\newtcolorbox{amdnote}{
    boxbasestyle,
    colback = cyan!5, colbacktitle = cyan!50!black, colframe = cyan!50!black,
    coltitle = white, title = {\faInfoCircle\quad \trAdmNote},
}
\newtcolorbox{amdtip}{
    boxbasestyle,
    colback = green!5!white, colbacktitle = green!40!black, colframe = green!40!black,
    coltitle = white, title = {\faLightbulb\quad \trAdmTip},
}
\newtcolorbox{amdimportant}{
    boxbasestyle,
    colback = violet!5, colbacktitle = violet!60!black, colframe = violet!60!black,
    coltitle = white, title = {\faExclamationCircle\quad \trAdmImportant},
}
\newtcolorbox{amdwarning}{
    boxbasestyle,
    colback = orange!5, colbacktitle = orange!70!black, colframe = orange!70!black,
    coltitle = white, title = {\faExclamationTriangle\quad \trAdmWarning},
}
\newtcolorbox{amdcaution}{
    boxbasestyle,
    colback = red!5, colbacktitle = red!70!black, colframe = red!70!black,
    coltitle = white, title = {\faBan\quad \trAdmCaution},
}

% ==============================================================
% 14. DRAFT WATERMARK
% ==============================================================
% main.tex sets \IsDraft to "true" to stamp every page "DRAFT" /
% "MINUTA" while the manuscript is a work in progress.
\providecommand{\IsDraft}{false}
\ifdefstring{\IsDraft}{true}{
    \usepackage{draftwatermark}
    \SetWatermarkText{\trDraft}
    \SetWatermarkScale{1}
    \SetWatermarkColor[gray]{0.8}
    \SetWatermarkAngle{30}
}{}

% ==============================================================
% 15. COVER
% ==============================================================
% Everything here reads macros defined in cover.tex, which
% build_cover.py generates from journal.toml, article.toml and
% authors.toml. Nothing on the cover is written by hand, which is
% why contents.tex can begin straight at the first section.
%
% \PrintMasthead fills the wide column: journal name, title,
% authors, affiliations, then both abstracts. \PrintSidebar fills
% the narrow one, in the order RBRH prints it.

% The author and affiliation blocks are long runs of proper nouns
% and institution names, which offer TeX very few legal break
% points. On the default settings the author list overflowed the
% column by about 10pt.
%
% Justified was tried first and does fix the overflow, with
% \tolerance raised and \emergencystretch added -- but it pays for
% it by hyphenating personal names across lines ("An-derson",
% "Fer-nando") and by leaving one very loose line. Breaking a
% surname is not acceptable in a masthead, so these blocks are set
% ragged right instead, with hyphenation switched off outright.
% The right edge is uneven; nothing overflows and no name breaks.
%
% To go back to justified, swap \raggedright and the two penalties
% for: \tolerance=4000 \emergencystretch=2em
%
% These are paragraph-breaking parameters, so they only take
% effect if the paragraph ends while they are still in scope --
% hence the \par inside each group below.
\newcommand{\RBRHnameblock}{%
  \raggedright
  \hyphenpenalty=10000
  \exhyphenpenalty=10000
}

% --- Wide column ---
\newcommand{\PrintMasthead}{%
  {\noindent \sffamily \Large\textbf{\JournalName}}\\
  {\noindent \large \textit{\JournalNameEn}}\\[3mm]
  \begin{center}
    \raggedright {\Large \sffamily \textbf{\color{AccentColor}\ArticleTitle}}\\[3mm]
  \end{center}
  {\RBRHnameblock\noindent\AuthorBlock\par}
  \vspace{2mm}
  % Affiliations in the same grey as the captions, so they read as
  % secondary to the author names without fading out like the
  % licence notice does.
  {\footnotesize\color{heavygray}\RBRHnameblock\noindent\AffiliationBlock\par}
  \vspace{2mm}
  \hrule
  \vspace{2mm}
  \PrintAbstracts
}

% Both languages are printed whatever \DocLanguage says: RBRH
% carries an English abstract and a Portuguese resumo regardless
% of the language the body is written in.
% The \small applies to the heading words only -- "Abstract",
% "Keywords:" and their Portuguese counterparts, which were
% sitting at full body size. The abstract text itself stays at
% \footnotesize.
\newcommand{\PrintAbstracts}{%
  \begin{spacing}{0.8}
  \small
  {\noindent \sffamily\textbf{\color{AccentColor}Abstract}\quad{\footnotesize\color{TextColor} \AbstractEn}}\\[1mm]
  {\noindent \sffamily\textbf{\color{AccentColor}Keywords:}\quad{\footnotesize\color{TextColor} \KeywordsEn}}\\[3mm]
  {\noindent \sffamily\textbf{\color{AccentColor}Resumo}\quad{\footnotesize\color{TextColor} \AbstractPt}}\\[1mm]
  {\noindent \sffamily\textbf{\color{AccentColor}Palavras-chave:}\quad{\footnotesize\color{TextColor} \KeywordsPt}}
  \end{spacing}
}

% --- Narrow column ---
% One helper per block, so the running order is readable and can
% be reshuffled by editing \PrintSidebar alone.
% The accent marks the structure of the column -- the heading of
% each block -- while the content under it stays in body colour.
\newcommand{\RBRHcoverblock}[2]{%
  {\color{AccentColor}\textbf{#1}}\\[1mm]
  #2
  \\[4mm]
}

\newcommand{\PrintSidebar}{%
  {
  \sffamily
  \scriptsize
  \raggedright
  % licence block at the foot overrides it with grey.
  \textbf{\MakeUppercase{\ArticleType}}\\
  \ifdefempty{\ArticleDOI}{}{DOI\;\ArticleDOI\\[4mm]}
  \RBRHcoverblock{\trHighlights:}{\ArticleHighlights}
  \RBRHcoverblock{\trCitation:}{\ArticleCitation}
  \RBRHcoverblock{\trCorrespondence:}{\CorrespondingBlock}
  \RBRHcoverblock{\trCompetingInterests:}{\CompetingInterests}
  \RBRHcoverblock{\trAcknowledgements:}{\Acknowledgements}
  \RBRHcoverblock{\trEditors:}{%
    \trEditorInChief:\\ \EditorInChief\\[1mm]
    \trAssociatedEditor:\\ \AssociatedEditor}
  \RBRHcoverblock{\trDataAvailability:}{\DataAvailability}
  \RBRHcoverblock{\trPublicationDates:}{%
    \trPublished\ \DatePublished\\
    \trIssue\ \ArticleIssue}
  \RBRHcoverblock{\trHistory:}{%
    \trReceived\ \DateReceived\\
    \trReviewed\ \DateReviewed\\
    \trAccepted\ \DateAccepted}
  % Push the licence to the foot of the column: everything above
  % sets at its natural height and this absorbs the slack, so the
  % notice sits on the baseline of the page however much or little
  % metadata precedes it.
  \vfill
  \PrintLicense
  }
}

% CC-BY notice, set in grey so it reads as a footer to the
% sidebar rather than as content. The badge is a vector PDF
% converted from assets/cc-by.svg.
\newcommand{\PrintLicense}{%
  {\color{midgray}
   \includegraphics[height=3.2mm]{\LicenseBadge}\\[1mm]
   \LicenseCopyright\\
   \LicenseStatement\par}
}

% ==============================================================
% 16. APPENDIX RENUMBERING
% ==============================================================
% Call \BeginAppendix in contents.tex to switch into lettered,
% restarted numbering (A, A.1, Figure A.1, ...) for supplementary
% material.
\newcommand{\BeginAppendix}{%
  \appendix
  \renewcommand{\thesection}{\Alph{section}}
  \renewcommand{\thesubsection}{\Alph{section}.\arabic{subsection}}
  \renewcommand{\thesubsubsection}{\Alph{section}.\arabic{subsection}.\arabic{subsubsection}}
  \setcounter{section}{0}
  \setcounter{equation}{0}
  \setcounter{figure}{0}
  \setcounter{table}{0}
  \setcounter{lstlisting}{0}
  \renewcommand{\theequation}{\Alph{section}.\arabic{equation}}
  \renewcommand{\thefigure}{\Alph{section}.\arabic{figure}}
  \renewcommand{\thetable}{\Alph{section}.\arabic{table}}
  \renewcommand{\thelstlisting}{\Alph{section}.\arabic{lstlisting}}
  \pagesectionvspace
  \auxlisttitle{\trAppendices}
  \label{sec:appendix}
  \vspace{10mm}
  \small
  % Headings one level down too. The body is already \small, which
  % left \Large section titles looking oversized against it.
  \titleformat{\section}
    {\bfseries\sffamily\large\color{HeadingColor}}
    {\makebox[8mm][l]{\thesection}}{2mm}{}
  \titleformat{\subsection}
    {\bfseries\sffamily\normalsize\color{HeadingColor}}
    {\makebox[8mm][l]{\thesubsection}}{2mm}{}
  \titleformat{\subsubsection}
    {\sffamily\small\color{HeadingColor}}
    {\makebox[10mm][l]{\thesubsubsection}}{4mm}{}
}

% ==============================================================
% 17. SHARED COUNTERS ACROSS PARACOL COLUMNS
% ==============================================================
% paracol makes every counter column-LOCAL by default -- its
% internal list of global counters ships as just {page}. Each
% column therefore keeps its own `figure`, `table` and `equation`.
% A figure written inside a column and one written in a
% \switchcolumn[1]*[...] spanning block are stepped in different
% counter states, so the spanning one restarts at 1 and the
% numbering appears to repeat. `figure` is shared instead: the two
% columns here are a metadata sidebar plus one article body, not
% parallel independent texts.
%
% Name the counter explicitly -- do NOT use \globalcounter*, and
% do NOT add `table` here. Both were tried and both break this
% document with "Missing \endcsname inserted" at the first
% \section (the starred form because it sweeps in every counter,
% including `table`). Bisected against this exact package set:
%   figure   -> safe, and is what fixes the numbering bug
%   equation -> safe, but unused; add only if an equation ever
%               goes inside a \switchcolumn[1]*[...] span
%   table    -> BREAKS the build, do not enable
% `footnote` is left alone too: paracol manages its scope itself.
\globalcounter{figure}

% ==============================================================
% 18. AUTHOR-FACING FIGURE AND TABLE MACROS
% ==============================================================
% So contents.tex never carries layout plumbing:
%
%   \FigureM[pos]{name}{caption}{label}        text width
%   \FigureS[pos][frac]{name}{caption}{label}  small -- frac of the
%                                              text width, default
%                                              0.67 -- and the only
%                                              size whose caption
%                                              sits beside the
%                                              image, spanning into
%                                              the sidebar column
%   \FigureL[pos]{name}{caption}{label}        full width, margin to
%                                              margin, matching the
%                                              width of the header
%                                              rule
%
%   \begin{TableM}[pos]{caption}{label} ... \end{TableM}  text width
%   \begin{TableL}[pos]{caption}{label} ... \end{TableL}  full width
%
% [pos] is the usual LaTeX float placement specifier and defaults
% to t (top). The table environments take the tabular material as
% their body and put the caption above it, per table convention;
% figures caption below.
%
% <name> is the bare file name: no folder, no extension. The file
% is resolved inside figures/, preferring .pdf, then .png, then
% .jpeg, then .jpg. A name that matches nothing degrades to a
% placeholder box rather than failing the build.
%
% M/S/L map onto real floats -- figure/table for M, and the
% double-column figure*/table* for S and L, which is what makes
% them span the sidebar. paracol supports both (it carries its own
% \pcol@floatplacement and spanning-float machinery), and the
% figure counter is shared across columns by section 19 above, so
% numbering follows source order. \label goes straight after
% \caption inside the float, where \@currentlabel is correct.

\newif\ifRBRHfigfound
\newcommand{\RBRHfigpath}{}

% Resolve figures/<name> against the supported extensions, in
% order of preference.
\newcommand{\RBRHfindfigure}[1]{%
  \IfFileExists{figures/#1.pdf}%
    {\RBRHfigfoundtrue \def\RBRHfigpath{figures/#1.pdf}}%
    {\IfFileExists{figures/#1.png}%
      {\RBRHfigfoundtrue \def\RBRHfigpath{figures/#1.png}}%
      {\IfFileExists{figures/#1.jpeg}%
        {\RBRHfigfoundtrue \def\RBRHfigpath{figures/#1.jpeg}}%
        {\IfFileExists{figures/#1.jpg}%
          {\RBRHfigfoundtrue \def\RBRHfigpath{figures/#1.jpg}}%
          {\RBRHfigfoundfalse \def\RBRHfigpath{}}}}}%
}

% Placeholder for a missing file: gray box in golden ratio
% (height = width / 1.618) carrying a red notice. \detokenize
% keeps names containing "_" printable.
\newlength{\RBRHfigwidth}
\newlength{\RBRHfigheight}
\newcommand{\RBRHmissingfigure}[2]{%
  \setlength{\RBRHfigwidth}{#1}%
  \setlength{\RBRHfigheight}{\dimexpr\RBRHfigwidth*1000/1618\relax}%
  \begingroup
    \setlength{\fboxsep}{0pt}%
    \colorbox{lightgray}{%
      \begin{minipage}[c][\RBRHfigheight][c]{\RBRHfigwidth}%
        \centering
        \textcolor{red}{\sffamily\small
          Figure \texttt{\detokenize{#2}} not found}%
      \end{minipage}%
    }%
  \endgroup
}

% The image at #1 wide, or the placeholder if it is missing.
\newcommand{\RBRHfigbody}[2]{%
  \ifRBRHfigfound
    \includegraphics[width=#1]{\RBRHfigpath}%
  \else
    \RBRHmissingfigure{#1}{#2}%
  \fi
}

% --- M: text width ---
\NewDocumentCommand{\FigureM}{O{t} m m m}{%
  \begin{figure}[#1]
    \centering
    \RBRHfindfigure{#2}%
    \RBRHfigbody{\linewidth}{#2}%
    \caption{#3}\label{#4}%
  \end{figure}%
}

% --- L: full width, margin to margin ---
\NewDocumentCommand{\FigureL}{O{t} m m m}{%
  \begin{figure*}[#1]
    \centering
    \RBRHfindfigure{#2}%
    \RBRHfigbody{\textwidth}{#2}%
    \caption{#3}\label{#4}%
  \end{figure*}%
}

% --- S: small, caption beside the image, spanning the sidebar ---
% The image width is captured BEFORE the float opens, while
% \linewidth still means the text column: inside a double-column
% float LaTeX widens it to \textwidth. The float body is typeset
% here at the call site, so the captured value is the right one.
\NewDocumentCommand{\FigureS}{O{t} O{0.67} m m m}{%
  \setlength{\smallfigfixedwidth}{#2\linewidth}%
  \begin{figure*}[#1]
    \RBRHfindfigure{#3}%
    \setlength{\sidecapimgwidth}{\smallfigfixedwidth}%
    \setlength{\sidecapgap}{6mm}%
    \setlength{\sidecapwidth}{\textwidth}%
    \addtolength{\sidecapwidth}{-\sidecapimgwidth}%
    \addtolength{\sidecapwidth}{-\sidecapgap}%
    \noindent
    % \vspace{0pt} as the first thing in each minipage is what makes
    % [t] align the TOPS. Without it, [t] aligns first baselines, and
    % a minipage holding just a graphic has its baseline at the
    % graphic's bottom -- which dropped the caption to the foot of
    % the image instead of level with its top. \abovecaptionskip is
    % zeroed for the same reason: it is meant to separate a figure
    % from a caption printed below it, not beside it.
    \begin{minipage}[t]{\sidecapwidth}
      \vspace{0pt}%
      \setlength{\abovecaptionskip}{0pt}%
      \caption{#4}\label{#5}%
    \end{minipage}%
    \hspace{\sidecapgap}%
    \begin{minipage}[t]{\sidecapimgwidth}
      \vspace{0pt}%
      \RBRHfigbody{\linewidth}{#3}%
    \end{minipage}%
  \end{figure*}%
}

% --- Tables: caption above the body, per table convention ---
% The body is wrapped in threeparttable, so the caption and any
% notes take the width of the tabular rather than running wider,
% and authors can add a notes block without extra scaffolding:
%
%   \begin{TableM}{caption}{label}
%     \begin{tabular}{ll} ... \end{tabular}
%     \begin{tablenotes}
%       \item Note text.
%     \end{tablenotes}
%   \end{TableM}
\NewDocumentEnvironment{TableM}{O{t} m m}{%
  \begin{table}[#1]
    \centering
    \begin{threeparttable}
      \caption{#2}\label{#3}%
}{%
    \end{threeparttable}
  \end{table}%
}

\NewDocumentEnvironment{TableL}{O{t} m m}{%
  \begin{table*}[#1]
    \centering
    \begin{threeparttable}
      \caption{#2}\label{#3}%
}{%
    \end{threeparttable}
  \end{table*}%
}
